% concatenated from entropy_arXiv_ver2.tex
\documentclass[12pt]{article}
\usepackage{amsmath,amssymb}
\usepackage{amscd}
\usepackage{comment}
\usepackage{color}
\usepackage{enumerate}
\usepackage{cases}
\usepackage{mathtools}
\usepackage[all]{xy}
\usepackage{amsthm}
\usepackage{pb-diagram}
\usepackage[top=25truemm,bottom=25truemm,left=15truemm,right=15truemm]{geometry}
\theoremstyle{definition}
\newtheorem{defi}{Definition}%[section]
\newtheorem{theo}[defi]{Theorem}
\newtheorem{lemm}[defi]{Lemma}
\newtheorem{prop}[defi]{Proposition}
\newtheorem{cor}[defi]{Corollary}
\newtheorem{ex}[defi]{Example}
\newtheorem{rem}[defi]{Remark}
\newtheorem{que}[defi]{Question}
\newtheorem{conj}[defi]{Conjecture}
\newtheorem{cond}[defi]{Condition}
\renewcommand{\proofname}{\it Proof.}
\newcommand{\sslash}{{/\mkern-6mu/}}
%%%%%%%%%%%%%%%%%%%
\def\ad{{\rm ad}}
\def\Aut{{\rm Aut}}
\def\det{{\rm det}}
\def\div{{\rm div}}
\def\pr{{\rm pr}}
\def\Spin{{\rm Spin}}
\def\Sp{{\rm Sp}}
\def\SO{{\rm SO}}
\def\SL{{\rm SL}}
\def\SU{{\rm SU}}
\def\U{{\rm U}}
\def\GL{{\rm GL}}
\def\O{{\rm O}}
\def\Ad{{\rm Ad}}
\def\End{{\rm End}}
\def\Exp{{\rm Exp}}
\def\skew{{\rm skew}}
\def\Exp{{\rm Exp}}
\def\Im{{\rm Im}}
\def\Gr{{\rm Gr}}
\def\Sym{{\rm Sym}}
\def\Hom{{\rm Hom}}
\def\coker{{\rm coker}}
\def\ind{{\rm ind}}
\def\det{{\rm det}}
\def\id{{\rm id}}
\def\Vol{{\rm Vol}}
\def\Re{{\rm Re}}
\def\Tr{{\rm Tr}}
\def\refe{\ast}
\def\ini{{\rm in}}
\def\vol{{\rm vol}}
\def\cyc{{\rm cyc}}
\def\max{{\rm max}}
\def\min{{\rm min}}

\def\H{{\mathbb H}}
\def\R{{\mathbb R}}
\def\Z{{\mathbb Z}}
\def\C{{\mathbb C}}
\def\N{{\mathbb N}}
\def\Q{{\mathbb Q}}
\def\K{{\mathbb K}}
\def\CP{{\mathbb C}{\mathbb P}}
\def\D{{\mathbb D}}

\def\sp{\mathfrak{sp}}
\def\so{\mathfrak{so}}
\def\spin{\mathfrak{spin}}
\def\g{\mathfrak{g}}
\def\t{\mathfrak{t}}
\def\h{\mathfrak{h}}
\def\s{\mathfrak{s}}
\def\k{\mathfrak{k}}
\def\u{\mathfrak{u}}

\def\A{{\cal A}}
\def\M{{\cal M}}
\def\F{{\cal F}}
%\def\D{{\cal D}}
\def\cH{{\cal H}}
\def\CC{{\cal C}}

\def\inum{{\sqrt{-1}}}

\providecommand{\keywords}[1]
{
\small 
\textbf{{Keywords---}} #1
}
\providecommand{\MSC}[1]
{
\small 
\textbf{{MSC---}} #1
}


\begin{document}

\title {Shannon entropy for harmonic metrics on cyclic Higgs bundles}
\author {Natsuo Miyatake}
\date{}
\maketitle
\begin{abstract} 
Let $X$ be a Riemann surface, $K_X \rightarrow X$ the canonical bundle, and $T_X\coloneqq K_X^{-1}\rightarrow X$ the dual bundle of the canonical bundle. For each integer $r \geq 2$, each $q \in H^0(K_X^r)$, and each choice of the square root $K_X^{1/2}$ of the canonical bundle, we canonically obtain a Higgs bundle, which is called a cyclic Higgs bundle. A diagonal harmonic metric $h = (h_1, \dots, h_r)$ on a cyclic Higgs bundle yields $r-1$-Hermitian metrics $H_1, \dots, H_{r-1}$ on $T_X\rightarrow X$, defined as $H_j \coloneqq h_j^{-1} \otimes h_{j+1}$ for each $j=1,\dots, r-1$, while $h_1$, $h_r$, and $q$ yield a degenerate Hermitian metric $H_r$ on $T_X \rightarrow X$. 
The $r$-differential $q$ induces a subharmonic weight function $\phi_q=\frac{1}{r}\log|q|^2$ on $K_X\rightarrow X$, and the diagonal harmonic metric depends solely on this weight function $\phi_q$. In the previous papers, the author introduced and studied the extension of harmonic metrics associated with arbitrary subharmonic weight function $\varphi$, which also constructs $r-1$-Hermitian metrics $H_1,\dots, H_{r-1}$ and a degenerate Hermitian metric $H_r$ on $T_X\rightarrow X$. In this paper, for each non-zero real parameter $\beta$, we introduce a function, which we call entropy, that quantifies the degree of mutual misalignment of the Hermitian metrics $H_1,\dots, H_r$. 
By extending the estimate established by Dai-Li and Li-Mochizuki to general subharmonic weight functions, we provide an upper bound and a lower bound for the entropy when $H_1,\dots, H_{r-1}$ are all complete and satisfy a condition concerning their approximation. Additionally, we show that the difference between the lower and upper bounds of entropy converges to a finite real number if and only if $\beta>-1$.


\end{abstract}

\MSC{30C15, 31A05, 53C07}

\keywords{Cyclic Higgs bundles, Harmonic metrics, Complete solution, Maximum principle, Zeros in holomorphic sections, Potential theory, Shannon entropy}

\section{Introduction}
As is well known, every compact connected Riemann surface possesses a constant Gaussian curvature K\"ahler metric. Among them, constant negative Gaussian curvature K\"ahler metrics and flat Gaussian curvature K\"ahler metrics are related to Hitchin's seminal example of Higgs bundles \cite{Hit1}. Hitchin's example is as follows: Let $X$ be a compact connected Riemann surface and $K_X\rightarrow X$ the canonical bundle. For each choice of the square root $K_X^{1/2}$ of the canonical bundle and each $q\in H^0(K_X^2)$, we can associate the following Higgs bundle $(\K_2,\Phi(q))\rightarrow X$:
\begin{align*}
&\K_2\coloneqq K_X^{1/2}\oplus K_X^{-1/2}, \\ 
&\Phi(q)\coloneqq 
\left(
\begin{array}{cc}
0 & q \\ 
1 &0 
\end{array}
\right)\in H^0(\End \K_2\otimes K_X).
\end{align*}
This example was introduced by Hitchin \cite[Section 11]{Hit1}, especially on a compact hyperbolic Riemann surface. Suppose that $X$ is a compact hyperbolic Riemann surface. Then $(\K_2,\Phi(q))\rightarrow X$ is always stable (cf. \cite{Hit1}), and thus there uniquely exists a solution $h$ to the Hitchin equation satisfying $\det(h)=1$. Moreover, from the symmetry of the Higgs field and the uniqueness of the solution to the Hitchin equation, the metric $h$ splits as $h=(h_1,h_1^{-1})$. Let $T_X\rightarrow X$ be the dual bundle of $K_X\rightarrow X$. The solution $h=(h_1,h_1^{-1})$ induces a metric $H_1\coloneqq h_1^{-2}$ on $T_X\rightarrow X$ and one degenerate metric (cf. \cite[Theorem (11.2)]{Hit1}) $H_2\coloneqq h_1^2\otimes h_q$ on $T_X\rightarrow X$, where $h_q$ denotes the degenerate Hermitian metric on $K_X^{-2}\rightarrow X$ induced by $q$, defined as follows:
\begin{align}
(h_q)_x(u,v)\coloneqq \langle q_x,u\rangle\overline{\langle q_x,v\rangle} \ \text{for $x\in X$, $u,v\in (K_X^{-2})_x$}. \label{uv}
\end{align}
If $q$ is identically zero, then the K\"ahler metric induced by $H_1$ is a constant negative Gaussian curvature metric (cf. \cite[Section 11]{Hit1}). If $X$ is an elliptic curve and $q$ is a non-zero section, then we can easily find a solution to the Hitchin equation; $h=(h_1, h_1^{-1})=(h_q^{-1/4}, h_q^{1/4})$ solves the Hitchin equation, where $h_q$ is a Hermitian metric on $K_X^{-2}\rightarrow X$ defined similarly as (\ref{uv}). In this case, if we define $H_1=h_1^{-2}$ and $H_2=h_1^2\otimes h_q$ in the same way as in the case where $X$ is hyperbolic, we obtain $H_1=H_2=h_q^{1/2}$, and the K\"ahler metrics induced by them are flat Gaussian curvature metrics. From the above, we can observe that the solution of the Hitchin equation for Hitchin's example contains both constant negative Gaussian curvature K\"ahler metrics and flat Gaussian curvature K\"ahler metrics. They are essentially equivalent to the harmonic metrics on the Hitchin's example for the two extremal cases; the case where $q$ is identically zero, and the case where it has no zeros at all. We can further observe that in the case where $H_1$ yields a hyperbolic metric, $H_1$ and $H_2$ do not coincide anywhere on $X$ since $H_2$ is identically zero, while in the case where $H_1$ is flat, $H_1$ and $H_2$ coincide everywhere. Hitchin's example can naturally be generalized to harmonic metrics on higher rank cyclic Higgs bundles (cf. \cite{Bar1, Hit2}). Similar observations as above can be made for harmonic metrics on the general cyclic Higgs bundles as will be explained in Section \ref{2}.

In previous papers \cite{Miy2, Miy3}, the author proposed a new direction of research for cyclic Higgs bundles that considers harmonic metrics even on cyclic ramified coverings and Hermitian metrics that solve the extended Hitchin equation with a ``degenerate" or ``collapsed" coefficient obtained as the limit of infinitely increasing covering order. In other words, this means that we consider a more general semipositive singular Hermitian metric $e^{-\varphi}h_\refe$ on the canonical bundle instead of the singular metric induced by $q$. The intention was to connect the potential theory, especially the theory of zero-point configurations of holomorphic functions and holomorphic sections which has been the subject of much research for a very long time (cf. \cite{BCHM1, FKMR1, Ran1, ST1, SZ1}), with the theory of harmonic metrics on cyclic Higgs bundles. In \cite{Miy2, Miy3}, the author introduced an elliptic equation that extends the Hitchin equation for the ordinary cyclic Higgs bundle to be adjunct to an arbitrary semipositive singular metric $e^{-\varphi} h_\refe$ on $K_X\rightarrow X$, and discussed some fundamental properties of the equation. As with the usual cyclic Higgs bundles, by taking the adjacent components of the solution $(h_1,\dots, h_r)$ to the elliptic equation, we obtain $r-1$-Hermitian metrics $H_1,\dots, H_{r-1}$ on $T_X\rightarrow X$ and $h_1, h_r$ and $e^{-\varphi}h_\refe$ yield one degenerate Hermitian metric $H_r$ on $T_X\rightarrow X$ (see Section \ref{2.2}).

In this paper, we introduce a new function which we call {\it entropy}. It is constructed from the metrics $H_1,\dots, H_r$. Our first main theorem concerns the uniform estimate of the entropy from above and below. The second theorem asserts the difference between the lower and upper bounds of entropy converges to a finite real number if and only if $\beta>-1$. More specifically, our main theorems are as follows: 
\begin{theo}[Theorem \ref{main theorem a}]\label{main theorem 1} {\it Let $e^{-\varphi}h_\refe$ be a semipositive singular Hermitian metric on $K_X\rightarrow X$ that is not identically $\infty$ and is non-flat. Suppose that there exists a complete solution $h=(h_1,\dots, h_r)$ to equation (\ref{phi}) in Section \ref{2.2} associated with $\varphi$ that satisfies Condition \ref{three} in Section \ref{2.2}. When $r=2,3$, suppose in addition that $e^{\varphi}h_\refe^{-1}$ belongs to $W^{1,2}_{loc}$.
Then, for any non-zero real number $\beta$, the entropy $S(r,\varphi, \beta)$ constructed from the complete solution $h$ satisfies the following uniform estimate:
\begin{align}
S_{r,\beta}\leq S(r,\varphi,\beta)(x)<\log r\ \text{for any $x\in X$}, \label{SrbSrp}
\end{align}
where $S_{r,\beta}$ is the entropy for the weight function which is identically $-\infty$. Moreover, the equality in the lower bound of $S(r,\varphi,\beta)(x)$ is achieved if and only if $r=2,3$ and $\varphi(x)=-\infty$. 
}
\end{theo}

\begin{theo}[Theorem \ref{main theorem b}]\label{main theorem 2} {\it The following holds:
\begin{align*}
\lim_{r\to\infty} (S_{r,\beta}-\log r)=
\begin{cases}
-\frac{2\beta d_\beta}{c_\beta}+\log(c_\beta) \ &\text{if $\beta>-1$} \\
-\infty \ &\text{if $\beta\leq -1$},
\end{cases}
\end{align*}
where $c_\beta$ and $d_\beta$ are defined as follows:
\begin{align*}
&c_\beta\coloneqq \int_0^1s^\beta(1-s)^\beta ds, \\
&d_\beta\coloneqq \int_0^1s^\beta(1-s)^\beta \log s\ ds.
\end{align*}
}
\end{theo}
Theorem 1 asserts that entropy can be uniformly evaluated from above and below using the entropy in two extreme cases: when the metric $e^{-\varphi}h_\refe$ is flat and when the weight $\varphi$ is identically minus infinity. The concept of entropy was strongly influenced by and motivated by the mutual domination theorem of harmonic metrics on cyclic Higgs bundles established by Dai-Li \cite{DL1, DL2} and Li-Mochizuki \cite{LM1} which asserts that the difference between $H_1,\dots, H_r$ is bounded from above and below by the difference of metrics of the two extreme cases (see Section \ref{2.1}). Most of the proof of Theorem \ref{main theorem 1} is devoted to extending the estimate established by Dai-Li and Li-Mochizuki to more general subharmonic weight functions.

Condition \ref{three} in Theorem \ref{main theorem 1} concerns the approximation of the complete solution. The author expects that for any non-trivial semipositive singular Hermitian metric $e^{-\varphi}h_\refe$ on $K_X\rightarrow X$, there always exists a unique complete solution to the extended Hitchin equation (\ref{phi}) in Section \ref{2.2} associated with $\varphi$, and that approximation of this complete solution by smooth complete solutions in the sense of Condition \ref{three} is always possible. However, the proof has not yet been completed (see Section \ref{2.2} and \cite{Miy5}).

The assumption in Theorem \ref{main theorem 1} that $e^\varphi h_\refe^{-1}$ belongs to $W^{1,2}_{loc}$ is made in order to apply the mean value inequality (see Section \ref{CYMMVI}) in the proof of Proposition \ref{r/2}. This assumption also seems unnecessary, but the proof is not yet complete without it. It is used in Proposition \ref{r/2} to assert that equality is never achieved in the inequality $H_0 \otimes H_1^{-1} \leq 1$. Therefore, if $r\geq 4$, estimate (\ref{SrbSrp}) holds even without this assumption.




In the definition of entropy mentioned above, $\beta$ is a non-zero real value based on the motif of inverse temperature in the canonical ensemble of statistical physics. In the paper \cite{Miy4}, which is a sequel to this paper, we will pursue the analogy with the canonical ensemble in more detail. 

The structure of the paper is organized as follows: In Section \ref{2}, we will explain some background. In Section \ref{CYMMVI}, we will give a brief review for the Cheng-Yau maximum principle and the mean value inequality, which we will use in the proof of Theorem \ref{main theorem 1}.
In Section \ref{4}, we will give the definition of entropy. In Section \ref{5}, we will give the statement of our main theorems. In Section \ref{6}, we will provide a proof of our main theorems. 

\medskip
\noindent
{\bf Acknowledgements.} 

\noindent
I am grateful to Ryoichi Kobayashi for questioning the direction of moving the parameter $r$, which controls the size of the target symmetric space of harmonic maps. 
I would like to express my deep gratitude to Qiongling Li for answering my many questions, for reading the first draft of this paper, and for pointing out an improvement in the upper bound of the entropy when the $r$-differential $q$ is bounded, with respect to the complete hyperbolic K\"ahler metric, by a constant depending only on the norm of $q$ measured by the hyperbolic metric. 
I would also like to extend my sincere thanks to Toshiaki Yachimura for his advice on the readability of this paper, for our many discussions, and for his invaluable help, especially in the definition of entropy. This work was supported by the Grant-in-Aid for Early Career Scientists (Grant Number: 24K16912).




\section{Backgrounds}\label{2}
\subsection{Real, diagonal, and complete Hermitian metrics}\label{rdc}
Let $X$ be a connected, possibly non-compact Riemann surface and $K_X\rightarrow X$ the canonical bundle. We choose a square root $K_X^{1/2}$ of the canonical bundle $K_X\rightarrow X$. Let $\K_r$ be a holomorphic vector bundle of rank $r$ defined as follows:
\begin{align*}
\K_r\coloneqq \bigoplus_{j=1}^r K_X^{\frac{r-(2j-1)}{2}} = K_X^{\frac{r-1}{2}} \oplus K_X^{\frac{r-3}{2}} \oplus \cdots \oplus K_X^{-\frac{r-3}{2}} \oplus K_X^{-\frac{r-1}{2}}.
\end{align*}
A Hermitian metric $h$ on $\K_r\rightarrow X$ is said to be {\it real} (cf. \cite{Hit2}) if the following bundle isomorphism $S$ is isometric with respect to $h$:
\begin{align*}
S\coloneqq \left(
\begin{array}{ccc}
& & 1 \\ 
& \reflectbox{$\ddots$} & \\ 
1 & &
\end{array}
\right): \K_r\rightarrow \K_r^{\vee},
\end{align*}
where $\K_r^{\vee}$ denotes the dual bundle of $\K_r$. Also, $h$ is called {\it diagonal} (cf. \cite{Bar1, Col1, DL2}) if it splits as $h=(h_1,\dots, h_r)$. Let $T_X\rightarrow X$ be the dual bundle of $K_X\rightarrow X$. By taking the difference between the adjacent components in a diagonal Hermitian metric $h=(h_1,\dots, h_r)$, we obtain $r-1$-Hermitian metrics $H_1,\dots, H_{r-1}$ on $T_X\rightarrow X$:
\begin{align*}
H_j\coloneqq h_j^{-1}\otimes h_{j+1}.
\end{align*}
A diagonal Hermitian metric $h$ is real if and only if $h_j=h_{r-j+1}$ for all $j=1,\dots, r$. This is also equivalent to the metrics $H_1,\dots, H_{r-1}$ satisfying $H_j=H_{r-j}$ for all $j=1,\dots, r-1$. We say that a diagonal Hermitian metric $h$ is {\it complete} (cf. \cite{LM1}) if the K\"ahler metrics induced by $H_1,\dots, H_{r-1}$ are all complete. 
\subsection{Harmonic metrics on cyclic Higgs bundles}\label{2.1}
Let $q$ be a holomorphic section of $K_X^r\rightarrow X$. Then we associate a holomorphic section $\Phi(q)\in H^0(\End \K_r\otimes K_X)$, which is called a Higgs field, as follows:
\begin{align*}
\Phi(q)\coloneqq 
\left(
\begin{array}{cccc}
0 & && q \\ 
1 & \ddots && \\ 
& \ddots & \ddots & \\ 
& & 1 & 0
\end{array}
\right).
\end{align*}
The pair $(\K_r,\Phi(q))$ is called a cyclic Higgs bundle (cf. \cite{Bar1, Hit1, Hit2}). The above cyclic Higgs bundle is an example of cyclotomic Higgs bundles introduced in \cite{Sim3}. For each Hermitian metric $h$ on $\K_r$, we associate a connection $D_q(h)$ defined as follows:
\begin{align*}
D_q(h)\coloneqq \nabla^h + \Phi(q) + \Phi(q)^{\ast h}.
\end{align*}
A Hermitian metric on $(\K_r,\Phi(q))\rightarrow X$ is called a {\it harmonic metric} if the connection $D_q(h)$ is a flat connection. A harmonic metric $h$ yields a harmonic map $\hat{h}: \widetilde{X}\rightarrow \SL(r,\C)/\SU(r)$, where $\widetilde{X}$ is the universal covering space. f $h$ is real, then $\hat{h}$ maps into $\SL(r,\R)/\SO(r)$. Suppose that $X$ is a compact Riemann surface of genus at least $2$. Then the Higgs bundle $(\K_r,\Phi(q))$ is stable for any $q\in H^0(K_X^r)$, and thus there uniquely exists a harmonic metric $h$ on $(\K_r,\Phi(q))$ such that $\det(h)=1$. From the uniqueness of the metric and the symmetry of the Higgs field (cf. \cite{Bar1, Hit2}), the harmonic metric $h$ is diagonal and real. A non-diagonal harmonic metric (resp. a non-real harmonic metric) can be constructed by considering the Dirichlet problem (cf. \cite{Don1, LM2}) for the Hitchin equation with a non-diagonal (resp. a non-real) Hermitian metric on the boundary. 
Let $h=(h_1,\dots, h_r)$ be a diagonal harmonic metric on a cyclic Higgs bundle $(\K_r,\Phi(q))\rightarrow X$. As mentioned in the previous subsection we obtain $r-1$-Hermitian metrics $H_1,\dots, H_{r-1}$ on $T_X\rightarrow X$ defined as
\begin{align*}
H_j\coloneqq h_j^{-1}\otimes h_{j+1} \ \text{for each $j=1,\dots,r -1$}.
\end{align*}
We also obtain a Hermitian metric $H_r$ on $T_X\rightarrow X$ which is degenerate at the zeros of $q$ as follows:
\begin{align*}
H_r\coloneqq h_r^{-1}\otimes h_1 \otimes h_q,
\end{align*}
where $h_q$ is a Hermitian metric on $K_X^{-r}\rightarrow X$ that is degenerate at the zeros of $q$ defined as follows:
\begin{align}
(h_q)_x(u,v)\coloneqq \langle q_x,u\rangle\overline{\langle q_x,v\rangle} \ \text{for $x\in X$, $u,v\in (K_X^{-r})_x$}. \label{q}
\end{align}
The connection $D_q(h)$ is flat if and only if the Hermitian metric $h$ satisfies the following elliptic equation, called the Hitchin equation:
\begin{align*}
F_h + [\Phi(q)\wedge\Phi(q)^{\ast h}] = 0.
\end{align*}
By using $H_1,\dots, H_r$, we can describe the Hitchin equation as follows: 
\begin{align}
\inum F_{h_j} + \vol(H_{j-1}) - \vol(H_j) = 0 \ \text{for $j=1,\dots, r-1$}, \label{F}
\end{align}
where for each $j=1,\dots, r$, $\vol(H_j)$ is the volume form of the metric $H_j$ which is locally described as $\vol(H_j)=\inum H_j(\frac{\partial}{\partial z}, \frac{\partial}{\partial z})\ dz\wedge d\bar{z}$, and $H_0$ is $H_r$. Equation (\ref{F}) is a kind of Toda equation (cf. \cite{AF1, Bar1, GH1, GL1, Moc0, Moc1}). By making a variable change from $h=(h_1,\dots, h_r)$ to $H_1,\dots, H_{r-1}$, the Hitchin equation becomes the following:
\begin{align}
\inum F_{H_j} + 2\vol(H_j) - \vol(H_{j-1}) - \vol(H_{j+1}) = 0 \ \text{for $j=1,\dots, r-1$.} \label{cyc}
\end{align}
In the case where $r=2$ and $q=0$, the Hitchin equation for the cyclic Higgs bundle is equivalent to the following equation:
\begin{align}
\inum F_H + 2\vol(H) = 0, \label{H}
\end{align}
where the solution $H$ to equation (\ref{H}) is a Hermitian metric $H$ on $T_X\rightarrow X$. The K\"ahler metric induced by $H$ is a constant negative Gaussian curvature K\"ahler metric (cf. \cite{Hit1}). 
We remark on the following two points:
\begin{itemize}
\item If $q$ has no zeros, then $H_1=\cdots=H_{r-1}=h_q^{1/r}$ is a solution to PDE (\ref{cyc}).
\item Suppose that $q=0$ and that there is a solution $H$ to equation (\ref{H}). We define positive constants $\lambda_1,\dots, \lambda_{r-1}$ as follows: 
\begin{align}
\lambda_j\coloneqq 2\sum_{k=1}^{r-1}(\Lambda_{r-1}^{-1})_{jk} = j(r-j), \label{lambda}
\end{align}
where $\Lambda_{r-1}^{-1}$ is the inverse matrix of the Cartan matrix of type $A_{r-1}$. Then $H_j=\lambda_j H\ (j=1,\dots, r-1)$ is a solution to PDE (\ref{cyc}).
\end{itemize}

Li-Mochizuki \cite{LM1} established the following result:
\begin{theo}[\cite{LM1}]{\it Suppose that $q$ is a non-zero holomorphic section unless $X$ is hyperbolic. Then there always exists a unique complete harmonic metric $h$ on $(\K_r,\Phi(q))\rightarrow X$ satisfying $\det(h)=1$. Moreover, the complete harmonic metric $h$ is real.
}
\end{theo}
In addition to \cite{LM1}, we refer the reader to \cite{DL3, DW1, GL1, LT1, LTW1, Li1, LM2, Moc0, Moc1, Nie1, Wan1, WA1} for the works related to complete harmonic metrics. Also, Dai-Li \cite{DL1, DL2} and Li-Mochizuki \cite{LM1} established the following:
\begin{theo}[\cite{DL1,DL2, LM1}]\label{estimate} 
{\it Let $q$ be an $r$-differential that is not identically zero and has at least one zero. Let $h$ be the diagonal complete harmonic metric on the cyclic Higgs bundle $(\K_r,\Phi(q))$ satisfying $\det(h)=1$. Then the following estimate holds:
\begin{align}
&\frac{\lambda_{j-1}}{\lambda_j}=\frac{(j-1)(r-j+1)}{j(r-j)} < H_{j-1}\otimes H_j^{-1} < 1 \ \text{for all $2\leq j\leq [r/2]$},\label{r/2}\\
& H_r\otimes H_1^{-1} < 1.
\end{align}
}
\end{theo}
The above theorem states that the difference between the adjacent metrics $H_{j-1}\otimes H_j^{-1}$ can be evaluated using the metrics for the above two extremal cases. For theorems on harmonic maps which are derived from the above estimate, see \cite{DL1, DL2, DL3, LM1}.

\begin{rem} When $r=2$ and $X$ is a compact hyperbolic Riemann surface, the inequality $H_2\otimes H_1^{-1}<1$ was shown by Hitchin in the proof of \cite[Theorem 11.2]{Hit1}.
As a consequence, it holds that $q+\omega_1+\omega_2+\bar{q}$ is positive-definite, where $q$ is considered to be a $(2,0)$-form, and $\omega_j \ (j=1,2)$ is a $(1,1)$-form induced by $H_j$ (see also \cite[Theorem 3.1]{Wol1}). 
\end{rem}
\begin{rem} Since \cite{Wol1}, how the accompanying harmonic metric depends on $t$ when the holomorphic $r$- differential $q$ is scale-transformed to $tq$ by the parameter $t$ has been well studied (cf. \cite{CL1, DL2, DW1, Lof1, LoTW1, Sak1, Wol1}), but there seems to be much less research investigating how the location of zeros affects the global properties of a harmonic metric. Note, however, that the local behavior of the harmonic metric on the cyclic Higgs bundle around the zeros of $q$, especially for $r=2,3$, has been well studied. See \cite[Section 1.8]{DW1} and the references cited therein. We also remark on \cite{Li1, LM1, LM2, Nie1} for studies on the relationship between the completeness of the metric induced by $q$ outside the zero sets and the uniqueness of the harmonic metric. There has also been a great deal of research into the zeros of holomorphic differentials, forgetting about harmonic metrics, especially in the case of $r=2,3$. See, for example, \cite{BS1, DW1, Kus1, Mas1} and the references cited therein.
\end{rem}
\subsection{More general subharmonic weight functions}\label{2.2}
For each $q\in H^0(K_X^r)$, the metric $h_q$ induces a singular metric (cf. \cite{GZ1}) $h_q^{-1/r}$ that diverges at the zeros of $q$. We can define a curvature for $h_q^{-1/r}$ in the sense of currents, which is semipositive and has support at the zero points of $q$. We fix a smooth metric $h_\refe$ and denote by $e^{-\phi_q}h_\refe$ the metric $h_q^{-1/r}$, where the weight function $\phi_q$ is defined as $\phi_q\coloneqq \frac{1}{r}\log|q|_{h_\refe}^2$. If we choose a local flat reference metric, then the weight function defines a local subharmonic function. Also, although we will not go into details here, there is a way to identify the weight function with a global psh function on the total space of the dual line bundle excluding zero points, see \cite[Section 2.1]{BB1}. In the sense described above, we call the function $\phi_q$ a {\it subharmonic weight function}, although this is a bit of an abuse of terminology. We consider a more general singular metric $e^{-\varphi}h_\refe$ with semipositive curvature, where the weight function $\varphi$ is a function that is locally a sum of a subharmonic function and a smooth function. For each $q_N\in H^0(K_X^N)$, let $\phi_{q_N}$ be defined as $\frac{1}{N}\log|q_N|_{h_\refe}^2$. Then $e^{-\phi_{q_N}}h_\refe$ is a singular metric with semipositive curvature. Moreover, any semipositive singular metric can be approximated, at least in the $L^1_{loc}$ sense, by a sequence $(e^{-\phi_{q_N}}h_\refe)_{N\in\N}$, where $q_N\in H^0(K_X^N)$ (cf. \cite{GZ1}). In \cite{Miy2, Miy3}, for each semipositive singular metric $e^{-\varphi}h_\refe$ and each $r\geq 2$, the following equation for a diagonal metric $h=(h_1,\dots, h_r)$ on $\K_r\rightarrow X$ was introduced:
\begin{align}
\inum F_{h_j} + \vol(H_{j-1}) - \vol(H_j) = 0 \ \text{for $j=1,\dots, r-1$}, \label{phi}
\end{align}
where $H_1,\dots, H_{r-1}$ are defined as $H_j\coloneqq h_j^{-1}\otimes h_{j+1}$ for each $j=1,\dots, r-1$, and $H_r=H_0$ is defined as follows:
\begin{align*}
H_r = H_0 &\coloneqq h_r^{-1}\otimes h_1 \otimes (e^{-\varphi}h_\refe)^{-r} \\ 
&= H_1^{-1}\otimes \cdots \otimes H_{r-1}^{-1} \otimes (e^{-\varphi}h_\refe)^{-r}.
\end{align*}
For the case where $\varphi=\phi_q=\frac{1}{r}\log|q|_{h_\refe}^2$, equation (\ref{phi}) becomes the Hitchin equation for the cyclic Higgs bundle $(\K_r,\Phi(q))$. For the case where $\varphi=\phi_{q_N}=\frac{1}{N}\log|q_N|_{h_\refe}^2$, equation (\ref{phi}) gives a harmonic metric on a ramified covering space of $X$ (see \cite[Section 2]{Miy2}). As mentioned above, more general $\varphi$ can be approximated by a sequence $(\phi_{q_N})_{N\in\N}$ at least in the $L^1_{loc}$-sense (cf. \cite{GZ1}), therefore equation (\ref{phi}) can be considered to be an equation obtained as the limit when the covering degree increases infinitely or the number of zeros increases infinitely. By making a variable change from $h=(h_1,\dots, h_r)$ to $H_1,\dots, H_{r-1}$, the equation becomes the following:
\begin{align}
\inum F_{H_j} + 2\vol(H_j) - \vol(H_{j-1}) - \vol(H_{j+1}) = 0 \ \text{for $j=1,\dots, r-1$.}\label{cyc2}
\end{align}
We remark on the following two points:
\begin{itemize}
\item If $e^{-\varphi}h_\refe$ is flat, then $H_1=\cdots=H_{r-1}=(e^{-\varphi}h_\refe)^{-1}$ is a solution to PDE (\ref{cyc2}).
\item Suppose that $\varphi=-\infty$. Then equation (\ref{cyc2}) is the same as equation (\ref{cyc}) for the case where $q=0$, and thus for a solution $H$ to equation (\ref{H}), $H_j=\lambda_j H\ (j=1,\dots, r-1)$ is a solution to PDE (\ref{cyc2}), where $\lambda_1,\dots, \lambda_{r-1}$ are constants defined in (\ref{lambda}).
\end{itemize}
Concerning the complete solutions to equation (\ref{phi}), the following holds:
\begin{theo}[\cite{Miy5}]\label{c1} {\it Suppose that $\varphi$ is not identically $-\infty$ unless $X$ is hyperbolic. Suppose also that $\varphi$ satisfies the following assumption:
\begin{enumerate}[$(\ast)$]
\item There exists a compact subset $K\subseteq X$ such that on $X\backslash K$, $e^{\varphi}$ is of class $C^2$.
\end{enumerate}
Then for any two complete solutions $h=(h_1,\dots, h_r)$ and $h^\prime=(h_1^\prime,\dots, h^\prime)$ associated with $\varphi$ satisfying $\det(h)=\det(h^\prime)=1$, we have $h=h^\prime$. 
}
\end{theo}

\begin{theo}[\cite{Miy5}]\label{c2} {\it On the unit disc $\D\coloneqq \{z\in\C\mid |z|<1\}$, for any $\varphi$, there exists a real complete solution $h=(h_1,\dots, h_r)$ associated with $\varphi$. For each $j=1,\dots, [r/2]$, $h_j$ is of class $C^{2j-1,\alpha}$ for any $\alpha\in(0,1)$. }
\end{theo}
\begin{theo}[\cite{Miy5}]\label{c3}{\it
Let $(e^{-\varphi_\epsilon}h_\refe)_{\epsilon>0}$ be a family of smooth semipositive Hermitian metrics on the canonical bundle, each of which is defined on a disc $\D_\epsilon \coloneqq \{z \in \C \mid |z| < 1 - \epsilon\}$, that satisfies the following property:
\begin{itemize}
\item For each $\epsilon>\epsilon^\prime>0$, $\varphi_{\epsilon^\prime}\leq \varphi_\epsilon$ and $(\varphi_\epsilon)_{0<\epsilon<1}$ converges to a function $\varphi$ on $\D$ that is locally a sum of a smooth function and a subharmonic function as $\epsilon\searrow 0$.
\end{itemize}
Then the corresponding family of smooth complete solutions $(h_\epsilon=(h_{1,\epsilon},\dots, h_{r,\epsilon}))_{\epsilon>0}$ monotonically converges to a complete solution $h=(h_1,\dots, h_r)$ associated with $\varphi$ as $\epsilon\searrow 0$.}
\end{theo}
The assumption $(\ast)$ in the uniqueness part of Theorem \ref{c1} is for using the maximum principle on open Riemann surfaces (cf. \cite[Section 3]{LM1}). The author expects that this assumption is unnecessary, but has not yet been able to prove the uniqueness without it. If we can drop the assumption $(\ast)$ in Theorem \ref{main theorem 1}, then, it follows from Theorem \ref{main theorem 2} that for every $\varphi$ on a hyperbolic surface, by lifting to the universal covering space and applying the uniqueness argument on the unit disc (cf. \cite[Proposition 5.7]{LM1}), there exists a complete solution associated with $\varphi$.

Theorem \ref{c3} on approximation for complete solutions is intended to approximately apply the various maximum principle techniques to complete solutions to which these techniques cannot be applied directly. Note that the approximation $(\varphi_\epsilon)_{0<\epsilon<1}$ of $\varphi$ in Theorem \ref{c3} always exists by the standard theory of mollification of subharmonic functions (cf. \cite{Ran1}). If Theorem \ref{c1} on the uniqueness of complete solutions is fully established, then for any $\varphi$ on a hyperbolic Riemann surface, the unique complete solution associated with $\varphi$ can be approximated by a smooth complete solution when lifted to the universal covering space. Similar results on parabolic Riemann surfaces have not yet been established but will be addressed in a subsequent paper. For the purpose of citation when stating the main theorem, we summarize the properties of approximations that we would like a complete solution to satisfy:

\begin{cond}\label{three} Let $h$ be a complete solution to equation (\ref{phi}) associated with $\varphi$. Let $\widetilde{X}$ denote the universal covering space of $X$. We choose the reference metric $h_\refe$ to be flat when pulled back to $\widetilde{X}$. Let $\widetilde{h}$ and $e^{-\widetilde{\varphi}}\widetilde{h}_\refe$ denote the pullbacks of $h$ and $e^{-\varphi}h_\refe$ to $\widetilde{X}$, respectively. Let $(\widetilde{\varphi}_\epsilon)_{\epsilon>0}$ be a mollification of $\widetilde{\varphi}$. We consider the following condition:
\begin{itemize}
\item There exists a family of complete solutions $(h_\epsilon)_{\epsilon>0}$, where each $h_\epsilon$ is associated with $\widetilde{\varphi}_\epsilon$, such that for every $x \in \widetilde{X}$, $h_\epsilon(x) \to \widetilde{h}(x)$ as $\epsilon \searrow 0$.
\end{itemize}
\end{cond}

\begin{rem} One of the author's initial motivation was to construct a theory that describes the variations of the ``random'' harmonic metrics $(h_N)_{N\in\N}$ associated with random sections $(q_N\in H^0(K_X^N))_{N\in\N}$ (cf. \cite{BCHM1, SZ1}). The author does not know if there is any research on harmonic metrics related to the above, but would like to remark on the following:
Hitchin equation for rank $2$ cyclic Higgs bundles has similarity to the PDE appeared in the Ginzburg-Landau model in physics (cf. \cite[Section 5]{DW1}, \cite[Section 11]{Hit1}, \cite{JT1},), and that the Ginzburg-Landau model has ties to the asymptotic behavior of empirical measures (cf. \cite{Ser1, Ser2}). Although indirect, this is the only known connection to the author between harmonic metrics and probability theory. 
The author is not aware of any work other than \cite{Miy2, Miy3} on harmonic metrics on cyclic ramified coverings 
considered in \cite{Miy2, Miy3}. 
\end{rem}


\subsection{Shannon entropy}\label{Shannon}
We provide a very quick review of Shannon entropy, which is the minimum required for use in this paper. The Shannon entropy is defined as follows:
\begin{defi}[\cite{Sh1}] 
Let $r\geq 2$ be a natural number and let $P_0,P_1,\dots, P_{r-1}$ be nonnegative numbers such that $P_0+P_1+\cdots+ P_{r-1}=1$. Then the Shannon entropy $S(P_0,P_1, \dots, P_{r-1})$ is defined as follows:
\begin{align*}
S(P_0, P_1,\dots, P_{r-1})\coloneqq -\sum_{j=0}^{r-1} P_j \log P_j.
\end{align*}
\end{defi}
It can be checked that we have the following estimate:
\begin{align*}
0\leq S(P_0,P_1,\dots, P_{r-1})\leq \log r.
\end{align*}
The minimum is attained if and only if $P_j=1$ for some $j=0,1,\dots,{r-1}$ and the others are 0, and the maximum is attained if and only if $P_0=P_1=\cdots =P_{r-1}=\frac{1}{r}$ (cf. \cite{Lei1, PW1, Sh1}). Shannon entropy quantifies the degree of bias of $P_0,P_1,\dots, P_{r-1}$. Shannon entropy is the basis for a very wide variety of entropy concepts, which we will not go into in depth in this paper. See, for example, \cite{KH1, Lei1, PW1} and the literature cited therein.

\section{The Cheng-Yau maximum principle and the mean value inequality}\label{CYMMVI}
This section briefly reviews the Cheng-Yau maximum principle \cite{CY1} (see also \cite[Section 3.1.2]{LM1}) and the mean value inequality \cite[Lemma 2.5]{CT1} (see also \cite[Section 2]{DL3}). The following theorem is known as the Cheng-Yau maximum principle, which played a crucial role in \cite{LM1}:
\begin{theo}[\cite{CY1}]\label{C-Ym}{\it Let $(M, g_M)$ be a complete Riemannian manifold with Ricci curvature bounded from below. Let $u$ be a real-valued $C^2$ function on $M$ satisfying $\Delta_{g_M} u \geq f(u)$, where $f:\R\rightarrow \R$ is a function, and where we denote by $\Delta_{g_M} = -d^\ast d$ the negative Laplacian. Suppose that there exists a continuous positive function $g:[a,\infty)\rightarrow \R_{>0}$ such that
\begin{enumerate}[(i)]
\item $g$ is non-decreasing;
\item $\liminf_{t\to\infty}\frac{f(t)}{g(t)}>0$;
\item $\int_a^\infty(\int_b^tg(\tau)d\tau)^{-1/2}dt<\infty$ for some $b\geq a$.
\end{enumerate}
Then the function $u$ is bounded above. Moreover, if $f$ is lower semicontinuous, then $f(\sup_M u)\leq 0$.
}
\end{theo}
The following theorem is known as the mean value inequality (see also \cite[Section 2]{DL3}):
\begin{theo}[\cite{CT1}]\label{MVI}{\it Let $(M,g_M)$ be a complete Riemannian manifold of dimension $n$. We consider the following differential inequality 
\begin{align}
\Delta_{g_M}u\leq c u \label{cuagain}
\end{align}
with some positive constant $c$. Let $x_0\in M$ be a point of $M$. We denote by $B_g(x_0,r)$ the open geodesic ball of radius $r$ centered at $x_0$, and by $\Vol(B(x_0, r))$ the volume of the geodesic ball. We fix a positive constant $R_0$ and suppose the following holds for $x_0$ and $R_0$:
\begin{enumerate}[(i)]
\item The Poincar\'e and the Sobolev inequalities hold for functions supported on $B_g(x_0, R_0)$ with constants $c_p$ and $c_s$;
\item There exists a positive constant $c_2$ such that $\Vol(B(x_0,r))\leq c_2r^n$ for all $r\leq R_0$.
\end{enumerate}
Then there exist positive constants $p_0$ and $C$ depending only on $n, c, c_2, c_p, c_s$ such that for any nonnegative $W^{1,2}$-function $u$ satisfying (\ref{cuagain}) on $B_g(x_0, R_0)$ and any $0<p<p_0$ the following inequality holds:
\begin{align*}
\inf_{x\in B_g(x_0,R_0/4)} u(x)\geq C\left(\int_{B_g(x_0,R_0/4)}u^pd\mu_g\right)^{1/p},
\end{align*}
where $d\mu_g$ is the volume measure. In particular, there exist constants $C>0$ and $0<p<1$ such that 
\begin{align*}
u(x_0)\geq C\left(\int_{B_g(x_0,R_0/4)}u^pd\mu_g\right)^{1/p}.
\end{align*}
}
\end{theo}
In \cite{CT1}, $u$ is assumed to be a $W^{1,2}$-function, but from reading the proof, we see that belonging to $W^{1,2}_{\mathrm{loc}}$ is sufficient. We will use Theorem \ref{C-Ym} and Theorem \ref{MVI} in Section \ref{6}. 
\section{Definition of entropy}\label{4}
Let $h=(h_1,\dots, h_r)$ be a solution to equation (\ref{phi}) and let $H_1,\dots, H_r$ be the Hermitian metrics constructed from $h$ and $e^{-\varphi}h_\refe$ as in Section \ref{2.2}. Then we define a function, which we call entropy, as follows:
\begin{defi}\label{entropy2} 
For each $j=0, 1,\dots, r-1$ and non-zero real number $\beta$, let $p_j(r,\beta, \varphi):X\rightarrow [0,1]$ be a nonnegative function defined as follows:
\begin{align*}
p_j(r,\beta, \varphi)\coloneqq \frac{\vol(H_j)^\beta}{\sum_{j=0}^{r-1}\vol(H_j)^\beta},
\end{align*}
where $\vol(H_j)^\beta/\sum_{j=0}^{r-1}\vol(H_j)^\beta$ is understood to be $(\vol(H_j)/\vol(H_\refe))^\beta/\sum_{j=0}^{r-1}(\vol(H_j)/\vol(H_\refe))^\beta$ for some reference metric $H_\refe$, which does not affect $p_j(r,\beta,\varphi)$. We call the following function entropy:
\begin{align*}
S(r,\beta, \varphi)\coloneqq -\sum_{j=0}^{r-1} p_j(r,\beta, \varphi)\log p_j(r,\beta, \varphi).
\end{align*}
\end{defi}
\begin{rem} The constant $\beta$ is modelled after the inverse temperature in the canonical ensemble in statistical mechanics (cf. \cite{LL1}). We will further pursue the analogy with the canonical ensemble in the subsequent paper \cite{Miy4}.
\end{rem}



\section{Main theorem}\label{5}
We introduce a notation. We denote by $S_{r,\beta}$ the entropy for the weight which is identically $-\infty$:
\begin{align*}
&S_{r,\beta}\coloneqq -\sum_{j=1}^{r-1} \hat{p}_{j,\beta}\log \hat{p}_{j,\beta}, \\
& \hat{p}_{j,\beta}\coloneqq \frac{\lambda_j^\beta}{\sum_{j=1}^{r-1}\lambda_j^\beta}.
\end{align*}
Our main theorems are as follows:
\begin{theo}\label{main theorem a} {\it Let $e^{-\varphi}h_\refe$ be a semipositive singular Hermitian metric on $K_X\rightarrow X$ that is not identically $\infty$ and is non-flat. Suppose that there exists a complete solution $h=(h_1,\dots, h_r)$ to equation (\ref{phi}) in Section \ref{2.2} associated with $\varphi$ that satisfies Condition \ref{three} in Section \ref{2.2}. Suppose that there exists a complete solution $h=(h_1,\dots, h_r)$ to equation (\ref{phi}) in Section \ref{2.2} associated with $\varphi$ that satisfies Condition \ref{three} in Section \ref{2.2}. When $r=2,3$, suppose in addition that $e^{\varphi}h_\refe^{-1}$ belongs to $W^{1,2}_{loc}$.
Then, for any non-zero real number $\beta$, the entropy $S(r,\varphi, \beta)$ constructed from the complete solution $h$ satisfies the following uniform estimate:
\begin{align*}
S_{r,\beta}\leq S(r,\varphi,\beta)(x)<\log r\ \text{for any $x\in X$},
\end{align*}
where $S_{r,\beta}$ is the entropy for the weight function which is identically $-\infty$. Moreover, the equality in the lower bound of $S(r,\varphi,\beta)(x)$ is achieved if and only if $r=2,3$ and $\varphi(x)=-\infty$. 
}
\end{theo}

\begin{theo}\label{main theorem b} {\it The following holds:
\begin{align*}
\lim_{r\to\infty} (S_{r,\beta}-\log r)=
\begin{cases}
-\frac{2\beta d_\beta}{c_\beta}+\log(c_\beta) \ &\text{if $\beta>-1$} \\
-\infty \ &\text{if $\beta\leq -1$},
\end{cases}
\end{align*}
where $c_\beta$ and $d_\beta$ are defined as follows:
\begin{align*}
&c_\beta\coloneqq \int_0^1s^\beta(1-s)^\beta ds, \\
&d_\beta\coloneqq \int_0^1s^\beta(1-s)^\beta \log s\ ds.
\end{align*}
}
\end{theo}
\begin{que}How does $S(r,\varphi,\beta)(x)$ behave along the variation of $r$ at each point $x$? For example, does the limit $\lim_{r\to\infty}(S(r,\varphi,\beta)(x)-\log r)$ exist?
\end{que}
\begin{rem} Consider the case where the weight function $\varphi$ equals $\phi_q=\frac{1}{r}\log|q|_{h_\refe}$ for some $q\in H^0(K_X^r)$. Suppose that $X$ is a hyperbolic Riemann surface and $\sup_X|q|_{g_X}$ is finite for the complete hyperbolic K\"ahler metric $g_X$ (for equivalent conditions, see \cite[Theorem 4.8]{DL3}). 
For simplicity, we also ignore the dependence of the entropy on $\beta$, for example by setting $\beta = 1$.
In this case, by applying Dai-Li's result \cite[Theorem 4.8]{DL3} (see also \cite{BH1, LT1, Wan1} for the works for the lower rank cases), entropy can be uniformly bounded from above by $\log r-\delta$, where $\delta$ is a positive constant which depends only on $\sup_X|q|_{g_X}$. 
The author would like to thank Qiongling Li for pointing this out when he sent her the first draft of the present paper. While the above improvement to the upper bound of the entropy will not be discussed in depth in this paper, it will be addressed in a subsequent work \cite{Miy4}.\end{rem} 


\begin{rem} The author is not aware of any research aimed at investigating the asymptotic behavior of the harmonic metric in some sense in the limit of the parameter $r$, which controls the size of the symmetric space to which the harmonic map goes, becoming infinitely large, but this seems like an interesting research direction. For research on the $r$-dependence of various quantities constructed from harmonic metrics, see, for example, \cite{DL1, DL2, Li0, LM3, Ma1}. In particular, by the estimate of Dai-Li \cite[Theorem 1.5]{DL2}, the sectional curvature of the image of the harmonic map associated with a cyclic Higgs bundle tends to vanish uniformly for all $q \in H^0(K_X^r)$ as the rank $r$ goes to infinity. 
\end{rem}



\section{Proof}\label{6}
\subsection{Proof of Theorem \ref{main theorem a}}
We first prove Theorem \ref{main theorem a}. The following proposition is an extension of the estimate established by Dai-Li \cite{DL1, DL2} and Li-Mochizuki \cite{LM1} to more general subharmonic weight functions:
\begin{prop}\label{r/2}
{\it Suppose that $\varphi$ is not identically $-\infty$ and that $e^{-\varphi}h_\refe$ is non-flat. Suppose that there exists a complete solution $h=(h_1,\dots, h_r)$ to equation (\ref{phi}) in Section {2.2} associated with $\varphi$ that satisfies Condition \ref{three} in Section \ref{2.2}. Let $H_1,\dots, H_r$ be Hermitian metrics on $K_X^{-1}\rightarrow X$ constructed from $h$ and $e^{-\varphi}h_\refe$ as in Section \ref{2.2}. Then the following estimate holds:
\begin{align}
&\frac{\lambda_{j-1}}{\lambda_j}=\frac{(j-1)(r-j+1)}{j(r-j)}<H_{j-1}\otimes H_j^{-1}<1 \ \text{for all $2\leq j\leq [r/2]$},\label{Gauss}\\
& H_r\otimes H_1^{-1}\leq 1. \label{Euler}
\end{align}
If, moreover, $e^{\varphi}h_\refe^{-1}$ belongs to $W^{1,2}_{loc}$, then $H_0\otimes H_1^{-1}<1$.
}
\end{prop}
We first prove the above Proposition \ref{r/2} for the case where $\varphi$ is smooth by following \cite[Theorem 4.4]{LM1}, and then we prove the assertion by using Condition \ref{three} and the mean value inequality (see Section \ref{CYMMVI}). We set $n\coloneqq [r/2]$ and $\delta\coloneqq r-2n$. Let $H_1,\dots, H_r$ be the Hermitian metrics on $K_X^{-1}\rightarrow X$ constructed by using a complete solution $h=(h_1,\dots, h_r)$ to equation (\ref{phi}) and $e^{-\varphi}h_\refe$. We define $\sigma_1,\dots, \sigma_n$ as follows:
\begin{align*}
\sigma_j\coloneqq \log(H_{j-1}\otimes H_j^{-1}) \ \text{for $j=1,2,\dots, n$.}
\end{align*}
Let $r\geq 4$. Then we also define $\sigma^\prime_1,\dots, \sigma_{n-1}^\prime$ as follows:
\begin{align*}
\sigma_j^\prime\coloneqq \log(H_j^{-1}\otimes H_{j+1}) \ \text{for $j=1,\dots, n-1$}.
\end{align*}
By direct calculation, we can check that the following holds:
\begin{align}
&\inum\partial\bar{\partial}\sigma_1\geq 3\vol(H_0)-4\vol(H_1)+\vol(H_2), \label{sigma} \\
&\inum \partial\bar{\partial}\sigma_j =3\vol(H_{j-1})-3\vol(H_j)-\vol(H_{j-2})+\vol(H_{j+1})\ \text{for $j=2,\dots, n-1$}, \label{sigma2}\\
&\inum \partial\bar{\partial} \sigma_n=(4-\delta)\vol(H_{n-1})-(3-\delta)\vol(H_n)-\vol(H_{n-2}), \label{sigma3}\\
&\inum \partial \bar{\partial} \sigma^\prime_j=-3\vol(H_j)+3\vol(H_{j+1})+\vol(H_{j-1})-\vol(H_{j+2})\ \text{for $j=1,\dots, n-2$}, \label{sigma4}\\
&\inum \partial \bar{\partial} \sigma^\prime_{n-1}=(3-\delta)\vol(H_n)-(4-\delta)\vol(H_{n-1})+\vol(H_{n-2}). \label{sigma5}
\end{align}
Note that we have used in inequality (\ref{sigma}) the assumption that the metric $e^{-\varphi}h_\refe$ is semipositive. We also note that when the regularity of $\varphi$ is not high enough, inequality (\ref{sigma}), equality (\ref{sigma2}) for $j=2$, and equality (\ref{sigma4}) hold only in the weak sense. %From (\ref{sigma}), we also have the following:
%\begin{align}
%\inum\partial \bar{\partial} e^{\sigma_1}\geq e^{\sigma_1}(3\vol(H_0)-4\vol(H_1)+\vol(H_2)) \label{esigma}.
%\end{align}
From (\ref{sigma}), (\ref{sigma2}), and (\ref{sigma3}), we obtain the following:
\begin{align}
&\inum\Lambda_{H_1}\partial \bar{\partial} \sigma_1\geq 3e^{\sigma_1}-4+e^{-\sigma_2}, \label{Hsigma}\\
&\inum \Lambda_{H_j}\partial\bar{\partial}\sigma_j =3e^{\sigma_j}-3-e^{\sigma_j+\sigma_{j-1}}+e^{-\sigma_{j+1}}\ \text{for $j=2,\dots, n-1$}, \label{Hsigma2} \\
&\inum \Lambda_{H_n}\partial\bar{\partial} \sigma_n=(4-\delta)e^{\sigma_n}-(3-\delta)-e^{\sigma_{n-1}+\sigma_n}. \label{Hsigma3}
\end{align}
From (\ref{sigma4}) and (\ref{sigma5}), we have
\begin{align}
&\inum \Lambda_{H_1}\partial \bar{\partial} \sigma^\prime_1\geq -3+3e^{\sigma_1^\prime}-e^{\sigma_1^\prime+\sigma_2^\prime}, \label{Hs}\\
&\inum \Lambda_{H_j}\partial \bar{\partial} \sigma^\prime_j=-3+3e^{\sigma_j^\prime}+e^{-\sigma_{j-1}^\prime}-e^{\sigma_j^\prime+\sigma_{j+1}^\prime}\ \text{$j=2,\dots, n-2$}, \label{Hsigma4}\\
&\inum \Lambda_{H_{n-1}}\partial \bar{\partial} \sigma^\prime_{n-1}=(3-\delta)e^{\sigma_{n-1}^\prime}-(4-\delta)+e^{-\sigma_{n-2}^\prime}. \label{Hsigma5}
\end{align}
We introduce the following notation:
\begin{align*}
&M_j\coloneqq \sup_Xe^{\sigma_j}\ \text{for $j=1,\dots, n$}, \\
&M_j^\prime\coloneqq \sup_Xe^{\sigma_j^\prime} \ \text{for $j=1,\dots, n-1$}, \\
&B_1\coloneqq 2(1-M_1^{-1}), \\
&B_j\coloneqq 2-M_j^{-1}-M_{j-1} \ \text{for $j=2,\dots, n$}, \\
&B_{n+1}\coloneqq (2-\delta)(1-M_n), \\
&B_0^\prime\coloneqq 2-M_1^\prime, \\
&B_j^\prime\coloneqq 2-M_j^{\prime-1}-M_{j+1}^\prime \ \text{for $j=1,\dots, n-2$}, \\
&B_{n-1}^\prime\coloneqq (2-\delta)(1-M_{n-1}^{\prime-1}).
\end{align*}
Following \cite[Section 4]{LM1}, with a slight variation of the argument, we first show the following:
\begin{lemm}\label{MMBB}{\it 
Suppose that $\varphi$ is smooth. Then the following holds:
\begin{align}
&M_j \leq \frac{4j}{2j+1}-\frac{2j-1}{2j+1}M_{j+1}^{-1}\ \text{for $j=1,\dots, n-1$}, \label{M}\\
&M_j B_j\leq B_{j+1}\ \text{for $j=1,\dots, n$}, \label{MB}\\
&M_{n-j}^\prime \leq \frac{4j-(2j-1)\delta}{2j+1-j\delta}-\frac{2j-1-(j-1)\delta}{2j+1-j\delta}M_{n-(j+1)}^{\prime-1} \ \text{for $j=1,\dots, n-2$}, \label{Mp}\\
&M_{n-j}^\prime B_{n-j}^\prime \leq B_{n-(j+1)}^\prime\ \text{for all $j=1,\dots, n-1$}. \label{MpBp}
\end{align}
}
\end{lemm}
\begin{proof}
For each Hermitian metric $H$ on $K_X^{-1}\rightarrow X$, we denote by $\Lambda_{H}$ the contraction operator of the K\"ahler metric induced by $H$. We first prove (\ref{M}) and (\ref{MB}) by induction on $j$. From (\ref{Hsigma}), we have 
\begin{align*}
\inum\Lambda_{H_1}\partial \bar{\partial} \sigma_1\geq 3e^{\sigma_1}-4+M_2^{-1}.
\end{align*}
From the Cheng-Yau maximum principle, we have
\begin{align*}
0\geq 3M_1-4+M_2^{-1}.
\end{align*}
This implies
\begin{align*}
M_1\leq \frac{4}{3}-\frac{1}{3}M_2^{-1} \ \text{and} \ M_1B_1\leq B_2.
\end{align*}
Therefore we obtain (\ref{M}) and (\ref{MB}) for the case where $j=1$.
Suppose that (\ref{M}) holds for some $j=k$, where $1\leq k<n-1$. We show that this implies inequality (\ref{MB}) for $j=k+1$. From (\ref{Hsigma2}), we have 
\begin{align*}
\inum \Lambda_{H_{k+1}}\partial\bar{\partial}\sigma_{k+1} \geq (3-M_k)e^{\sigma_{k+1}}-3+M_{k+2}^{-1}.
\end{align*}
Since we have (\ref{M}) for the case where $j=k$, it holds that $3-M_k>0$. Then from the Cheng-Yau maximum principle, we have
\begin{align*}
0\geq (3-M_k)M_{k+1}-3+M_{k+2}^{-1}.
\end{align*}
This implies 
\begin{align*}
M_{k+1}B_{k+1} \leq B_{k+2}.
\end{align*}
Suppose that (\ref{M}) holds for $j=n-1$. We show that this implies inequality (\ref{MB}) for $j=n$. From (\ref{Hsigma3}), we have
\begin{align*}
\inum \Lambda_{H_n}\partial\bar{\partial} \sigma_n\geq (4-\delta-M_{n-1})e^{\sigma_n}-(3-\delta). 
\end{align*}
Since we have (\ref{M}) for the case where $j=n-1$, it holds that $4-\delta-M_{n-1}>0$. Then from the Cheng-Yau maximum principle, we have 
\begin{align*}
0\geq (4-\delta-M_{n-1})M_n-(3-\delta).
\end{align*}
This implies (\ref{MB}) for the case where $j=n$. Suppose that (\ref{MB}) holds for some $j=k$, where $2\leq k\leq n-1$ and that inequality (\ref{M}) holds for $j=k-1$. We show that this implies (\ref{M}) holds for $j=k$, which completes the proof of (\ref{M}) and (\ref{MB}). From (\ref{MB}) for $j=k$, we have
\begin{align}
M_k(3-M_{k-1})\leq 3-M_{k+1}^{-1}. \label{mk}
\end{align}
By combining (\ref{mk}) with (\ref{M}) for $j=k-1$, we have
\begin{align*}
M_k\left(3-\frac{4k-4}{2k-1}+\frac{2k-3}{2k-1}M_k^{-1}\right)\leq 3-M_{k+1}^{-1}.
\end{align*}
This implies
\begin{align*}
M_k \leq \frac{4k}{2k+1}-\frac{2k-1}{2k+1}M_{k+1}^{-1}.
\end{align*} 
By induction on $j$, we obtain (\ref{M}) and (\ref{MB}) from the above. We next prove (\ref{Mp}) and (\ref{MpBp}). From (\ref{Hsigma5}), we have 
\begin{align*}
\inum \Lambda_{H_{n-1}}\partial \bar{\partial} \sigma^\prime_{n-1}\geq (3-\delta)e^{\sigma_{n-1}^\prime}-(4-\delta)+M_{n-2}^{\prime-1}. 
\end{align*}
Then from the Cheng-Yau maximum principle, we have
\begin{align*}
0\geq (3-\delta)M_{n-1}^\prime-(4-\delta)+M_{n-2}^{\prime-1}.
\end{align*}
This implies
\begin{align*}
M_{n-1}^\prime\leq \frac{4-\delta}{3-\delta}-\frac{M_{n-2}^{\prime-1}}{3-\delta} \ \text{and} \
M_{n-1}^\prime B_{n-1}^\prime\leq B_{n-2}^\prime,
\end{align*}
and thus we have proved (\ref{Mp}) and (\ref{MpBp}) for the case where $j=1$. Suppose that (\ref{Mp}) holds for some $j=k$, where $1\leq k<n-2$. We show that this implies (\ref{MpBp}) for $j=k+1$. From (\ref{Hsigma4}), we have
\begin{align*}
\inum \Lambda_{H_{n-(k+1)}}\partial \bar{\partial} \sigma^\prime_{n-(k+1)}\geq -3+3e^{\sigma_{n-(k+1)}^\prime}+M_{n-(k+2)}^{\prime-1}-M_{n-k}^\prime e^{\sigma_{n-(k+1)}^\prime}
\end{align*}
Since we have (\ref{Mp}) for $j=k$, it holds that $3-M_{n-k}^\prime>0$. Then by the Cheng-Yau maximum principle, we have
\begin{align*}
0\geq -3+3M_{n-(k+1)}^\prime+M_{n-(k+2)}^{\prime-1}-M_{n-k}^\prime M_{n-(k+1)}^\prime.
\end{align*}
This implies $M_{n-(k+1)}^\prime B_{n-(k+1)}^\prime\leq B_{n-(k+2)}^\prime$. Suppose that (\ref{Mp}) holds for $j=n-2$. We show that this implies (\ref{MpBp}) for $j=n-1$. From (\ref{Hs}), we have 
\begin{align*}
\inum \Lambda_{H_1}\partial \bar{\partial} \sigma^\prime_1\geq -3+3e^{\sigma_1^\prime}-e^{\sigma_1^\prime}M_2^\prime.
\end{align*}
Since we have (\ref{Mp}) for $j=n-2$, it holds that $3-M_2^\prime>0$. Then from the Cheng-Yau maximum principle, we have
\begin{align*}
0\geq -3+3M_1^\prime -M_1^\prime M_2^\prime.
\end{align*}
This implies $M_1^\prime B_1^\prime\leq B_0^\prime$. Suppose that (\ref{MpBp}) holds for some $j=k$, where $2\leq k\leq n-2$, and (\ref{Mp}) holds for $j=k-1$. We show that this implies (\ref{Mp}) for $j=k$, which completes the proof of (\ref{Mp}) and (\ref{MpBp}). From (\ref{MpBp}) for $j=k$, we have
\begin{align*}
M_{n-k}^\prime(3-M_{n-(k-1)}^\prime)\leq {3-M_{n-(k+1)}^{\prime-1}}.
\end{align*}
Since we have (\ref{Mp}) for $j=k-1$, it holds that
\begin{align*}
M_{n-k}^\prime\left(3-\frac{4(k-1)-(2k-3)\delta}{2k-1-(k-1)\delta}+\frac{2k-3-(k-2)\delta}{2k-1-(k-1)\delta}M_{n-k}^{\prime-1}\right)\leq {3-M_{n-(k+1)}^{\prime-1}}.
\end{align*}
This implies
\begin{align*}
M_{n-k}^\prime \leq \frac{4k-(2k-1)\delta}{2k+1-k\delta}-\frac{2k-1-(k-1)\delta}{2k+1-k\delta}M_{n-(k+1)}^{\prime-1}.
\end{align*}
Therefore we obtain (\ref{Mp}) for $j=k$. This implies (\ref{Mp}) and (\ref{MpBp}). 
\end{proof}

\begin{lemm}
{\it Suppose that $e^\varphi$ is smooth. Then the following holds:
\begin{align}
&M_j-1\leq \frac{2j-1}{2} B_{j+1},\ \text{for $j=1,\dots, n$}\label{mj1}\\
&M_{n-j}^\prime-1\leq \frac{2j-1-(j-1)\delta}{2-\delta}B_{n-(j+1)}^\prime\ \text{for $j=1,\dots, n-1$}.\label{mjp1}
\end{align}
}
\end{lemm}
\begin{proof}
We provide proofs of (\ref{mj1}) and (\ref{mjp1}) inductively with respect to $j$, which are identical to those in \cite[Proof of Claim 4.2 (3)]{LM1}. We first consider the case where $j=1$. From the definition of $B_1$ and $B_1^\prime$, we have the following:
\begin{align*}
M_1-1&=M_1(1-M_1^{-1}) \\
&=\frac{1}{2}M_1B_1 \\
&\leq \frac{1}{2}B_2, \\
M_{n-1}^\prime-1&=M_{n-1}^\prime(1-M_{n-1}^{\prime-1}) \\
&=\frac{1}{2-\delta}M_{n-1}^\prime B_{n-1}^\prime \\
&\leq \frac{1}{2-\delta} B_{n-2}^\prime,
\end{align*}
and thus we have (\ref{mj1}) and (\ref{mjp1}) for the case where $j=1$.
We next suppose that (\ref{mj1}) holds for $j=k$, where $1\leq k\leq n-1$, and that (\ref{mjp1}) holds for $j=k$, where $1\leq k\leq n-2$. We show that (\ref{mj1}) holds for $j=k+1$. By the definition of $B_{k+1}$, we have 
\begin{align*}
M_{k+1}-1&=M_{k+1}(1-M_{k+1}^{-1}) \\
&=M_{k+1}(B_{k+1}-1+M_k) \\
&\leq M_{k+1}(B_{k+1}+\frac{2k-1}{2}B_{k+1}) \\
&=\frac{2k+1}{2}M_{k+1}B_{k+1} \\
&\leq \frac{2k+1}{2}B_{k+2}, \\
M_{n-(k+1)}^\prime -1&=M_{n-(k+1)}^\prime(1-M_{n-(k+1)}^{\prime-1}) \\
&=M_{n-(k+1)}^\prime(B_{n-(k+1)}^\prime-1+M_{n-k}^\prime) \\
&\leq M_{n-(k+1)}^\prime(B_{n-(k+1)}^\prime+\frac{2k-1-(k-1)\delta}{2-\delta}B_{n-(k+1)}) \\
&=\frac{2k+1-k\delta}{2-\delta}M_{n-(k+1)}^\prime B_{n-(k+1)}^\prime \\
&\leq \frac{2k+1-k\delta}{2-\delta}B_{n-(k+2)}^\prime.
\end{align*}
This implies (\ref{mj1}) and (\ref{mjp1}) for $j=k+1$, and thus we have completed the proofs of (\ref{mj1}) and (\ref{mjp1}). 
\end{proof}
\begin{lemm}\label{26}{\it Suppose that $\varphi$ is smooth. Then the following holds:
\begin{align}
&M_j\leq 1\ \text{for $j=1,\dots, n$}, \label{MJ1} \\
&M_j^\prime\leq \frac{\lambda_{j+1}}{\lambda_j} \ \text{for $j=1,\dots, n-1$}. \label{lambdajj}
\end{align}
}
\end{lemm}
\begin{proof} The proof is identical to that in \cite[pp.23–24]{LM1}. We first prove (\ref{MJ1}). From (\ref{mj1}) for $j=n$, we have 
\begin{align*}
M_n-1\leq \frac{2n-1}{2}(2-\delta)(1-M_n).
\end{align*}
This implies $M_n\leq 1$. Suppose that (\ref{MJ1}) holds for some $j=k$, where $2\leq k\leq n$. We show that this implies (\ref{MJ1}) for $j=k-1$. From (\ref{mj1}) for $j=k$, we have 
\begin{align*}
M_{k-1}-1&\leq \frac{2k-3}{2}(2-M_k^{-1}-M_{k-1}) \\
&\leq \frac{2k-3}{2}(1-M_{k-1}).
\end{align*}
This implies $M_{k-1}\leq 1$, and thus we have proved (\ref{MJ1}). We next prove (\ref{lambdajj}). We introduce the following symbols:
\begin{align*}
d_j^\prime&\coloneqq \frac{\lambda_{j+1}}{\lambda_j}\ \text{for $j=1,\dots, n-1$}, \\
D_0^\prime&\coloneqq 2-d_1^\prime, \\
D_j^\prime&\coloneqq 2-d_j^{\prime-1}-d_{j+1}^\prime \ \text{for $j=1,\dots, n-2$}, \\
D_{n-1}^\prime&\coloneqq (2-\delta)(1-d_{n-1}^{\prime-1}).
\end{align*}
Then from the definition of $d_1^\prime,\dots, d_{n-1}^\prime$, the following holds:
\begin{align}
&0=-3+3d_1^\prime-d_1^\prime d_2^\prime, \\
&0=-3+3d_j^\prime+d_{j-1}^{\prime-1}-d_j^\prime d_{j+1}^\prime\ \text{$j=2,\dots, n-2$}, \\
&0=(3-\delta)d_{n-1}^\prime-(4-\delta)+d_{n-2}^{\prime-1}. 
\end{align}
Therefore from the definition of $D_1^\prime,\dots, D_{n-1}^\prime$ and the proof of Lemma \ref{MMBB}, we can check that the following holds:
\begin{align}
d_{n-j}^\prime-1=\frac{2j-1-(j-1)\delta}{2-\delta}D_{n-(j+1)}^\prime\ \text{for $j=1,\dots, n-1$}. \label{djp1}
\end{align}
We provide a proof of (\ref{lambdajj}). From (\ref{mjp1}) and (\ref{djp1}) for $j=n-1$, we have
\begin{align*}
M_1^\prime-1&\leq \frac{2n-3-(n-2)\delta}{2-\delta} (2-M_1^\prime), \\
d_1^\prime-1&=\frac{2n-3-(n-2)\delta}{2-\delta} (2-d_1^\prime).
\end{align*}
This implies $M_1^\prime\leq d_1^\prime$. Suppose that (\ref{lambdajj}) holds for some $j=k$, where $1\leq k\leq n-2$. We show that this implies (\ref{lambdajj}) for $j=k+1$. From (\ref{mjp1}) and (\ref{djp1}) for $j=n-(k+1)$, we have 
\begin{align*}
M_{k+1}^\prime-1&\leq \frac{2(n-k)-3-(n-k-2)\delta}{2-\delta}(2-M_k^{\prime-1}-M_{k+1}^\prime) \\
&\leq \frac{2(n-k)-3-(n-k-2)\delta}{2-\delta}(2-d_k^{\prime-1}-M_{k+1}^\prime), \\
d_{k+1}^\prime-1&=\frac{2(n-k)-3-(n-k-2)\delta}{2-\delta}(2-d_k^{\prime-1}-d_{k+1}^\prime).
\end{align*}
This implies $M_{k+1}^\prime\leq d_{k+1}^\prime$, which completes the proof of (\ref{lambdajj}).
\end{proof}
\begin{lemm}{\it Suppose that $\varphi$ is not identically $-\infty$ and that $e^{-\varphi}h_\refe$ is non-flat. Suppose also that $\varphi$ is smooth. Then the following estimate holds:
\begin{align*}
&\frac{\lambda_{j-1}}{\lambda_j}=\frac{(j-1)(r-j+1)}{j(r-j)}<H_{j-1}\otimes H_j^{-1}<1 \ \text{for all $2\leq j\leq n$},\\
& H_r\otimes H_1^{-1}<1.
\end{align*}
}
\end{lemm}
\begin{proof} 
We first consider $\sigma_1,\dots, \sigma_n$. From (\ref{MJ1}), (\ref{Hsigma}), (\ref{Hsigma2}), and (\ref{Hsigma3}), we have
\begin{align}
&\inum\Lambda_{H_1}\partial \bar{\partial} \sigma_1\geq 3e^{\sigma_1}-3, \label{sigmad1}\\
&\inum \Lambda_{H_j}\partial\bar{\partial}\sigma_j \geq 2e^{\sigma_j}-2 \ \text{for $j=2,\dots, n-1$}. \label{sigmad2} \\
&\inum \Lambda_{H_n}\partial\bar{\partial} \sigma_n\geq (3-\delta)(e^{\sigma_n}-1). \label{sigmad3}
\end{align}
From (\ref{MJ1}), (\ref{sigmad1}), (\ref{sigmad2}), (\ref{sigmad3}), and the strong maximum principle, we see that for each $j=1,\dots, n$, either of the following holds: $e^{\sigma_j}<1$ or $e^{\sigma_j}=1$, where the latter equality means that $e^{\sigma_j}$ is identically 1 on $X$. If $e^{\sigma_1}=1$, then from (\ref{Hsigma}), we have $e^{\sigma_2}=1$. If $e^{\sigma_j}=1$ for some $j=2,\dots, n-1$, then from (\ref{Hsigma2}), we have $e^{\sigma_{j-1}}=e^{\sigma_{j+1}}=1$. If $e^{\sigma_n}=1$, then from (\ref{Hsigma3}), we have $e^{\sigma_{n-1}}=1$. Note that if $e^{\sigma_1}=\cdots =e^{\sigma_n}=1$, then $e^{-\varphi}h_\refe$ is flat. Therefore if $e^{-\varphi}h_\refe$ is non-flat, then we have $e^{\sigma_j}<1$ for all $j=1,\dots, n$. We next consider $\sigma_1^\prime,\dots, \sigma_{n-1}^\prime$. From (\ref{lambdajj}), (\ref{Hs}), (\ref{Hsigma4}), and (\ref{Hsigma5}), we have
\begin{align}
&\inum \Lambda_{H_1}\partial \bar{\partial} \sigma^\prime_1\geq -3+3e^{\sigma_1^\prime}-d_2^\prime e^{\sigma_1^\prime}, \label{sigmadp1} \\
&\inum \Lambda_{H_j}\partial \bar{\partial} \sigma^\prime_j=-3+3e^{\sigma_j^\prime}+d_{j-1}^{\prime-1}-d_{j+1}^\prime e^{\sigma_j^\prime}\ \text{$j=2,\dots, n-2$}, \label{sigmadp2}\\
&\inum \Lambda_{H_{n-1}}\partial \bar{\partial} \sigma^\prime_{n-1}=(3-\delta)e^{\sigma_{n-1}^\prime}-(4-\delta)+d_{n-2}^{\prime-1}. \label{sigmadp3}
\end{align}
From (\ref{lambdajj}), (\ref{sigmadp1}), (\ref{sigmadp2}), and (\ref{sigmadp3}), and the strong maximum principle, we see that for each $j=1,\dots, n-1$, either of the following holds: $e^{\sigma_j^\prime}<d_j^\prime$ or $e^{\sigma_j^\prime}=d_j^\prime$. If $e^{\sigma_1^\prime}=d_1^\prime$, then from (\ref{Hs}), we have $e^{\sigma_2^\prime}=d_2^\prime$. If $e^{\sigma_j^\prime}=d_j^\prime$ for some $j=2,\dots, n-2$, then from (\ref{Hsigma4}), we have $e^{\sigma_{j-1}}=d_{j-1}$ and $e^{\sigma_{j+1}^\prime}=d_{j+1}^\prime$. If $e^{\sigma_{n-1}^\prime}=d_{n-1}^\prime$, then from (\ref{Hsigma5}), we have $e^{\sigma_{n-2}^\prime}=d_{n-2}^\prime$. Note that if $e^{\sigma_j^\prime}=d_j^\prime$ for all $j=1,\dots, n-1$, then from (\ref{sigma4}) for $j=1$ we see that $\varphi$ is identically $-\infty$. Therefore if $\varphi$ is not identically $-\infty$, then we have $e^{\sigma_j^\prime}<d_j^\prime$ for all $j=1,\dots, n-1$. This implies the claim. 
\end{proof}

Then we prove Proposition \ref{r/2}.
\begin{proof}[Proof of Proposition \ref{r/2}] Let $h=(h_1,\dots, h_r)$ be a complete solution to equation (\ref{phi}) in Section {2.2} associated with $\varphi$ that satisfies Condition \ref{three} in Section \ref{2.2}. Since it is sufficient to lift $X$ to the universal covering space and show the proposition, we will assume from the beginning that $h$ has an approximation $(h_\epsilon)_{\epsilon>0}$ in Condition \ref{three}. We denote by $H_{j,\epsilon} \ (j=1,\dots, r-1)$ the complete metrics associated with $e^{-\varphi_\epsilon}h_\ast$ and by $H_{r,\epsilon}$ the metric $ H_{1,\epsilon}^{-1}\otimes \cdots \otimes H_{r-1,\epsilon}^{-1} \otimes (e^{-\varphi_\epsilon}h_\refe)^{-r}$. Moreover, we set $\sigma_{j,\epsilon}\coloneqq \log(H_{j-1,\epsilon}\otimes H_{j,\epsilon}^{-1})$ for $j=1,2,\dots, n$ and $\sigma_{j,\epsilon}^\prime\coloneqq \log(H_{j,\epsilon}^{-1}\otimes H_{j+1,\epsilon})$ for $j=1,\dots, n-1$. We have proved (\ref{Gauss}) and (\ref{Euler}) for the case where $\varphi$ is smooth. From (\ref{Gauss}) and (\ref{Euler}) for $e^{-\varphi_\epsilon}h_\refe$, we have
\begin{align}
&\frac{\lambda_{j-1}}{\lambda_j}=\frac{(j-1)(r-j+1)}{j(r-j)}<H_{j-1,\epsilon}\otimes H_{j,\epsilon}^{-1}<1 \ \text{for all $2\leq j\leq n$},\label{Gausse}\\
& H_{r,\epsilon}\otimes H_{1,\epsilon}^{-1}<1. \label{Eulere}
\end{align}
By taking the limit $\epsilon\searrow 0$ of (\ref{Gausse}) and (\ref{Eulere}) at each point of $X$, we have 
\begin{align}
&\frac{\lambda_{j-1}}{\lambda_j}=\frac{(j-1)(r-j+1)}{j(r-j)}\leq H_{j-1}\otimes H_j^{-1}\leq 1 \ \text{for all $2\leq j\leq n$},\label{Gauss0}\\
& H_r\otimes H_1^{-1}\leq 1. \label{Euler0}
\end{align}
Our goal is to show that $e^{\sigma_j} < 1$ for each $j = 2, \dots, n$, that $e^{\sigma_j^\prime} < d_j^\prime = \frac{\lambda_{j+1}}{\lambda_j}$ for each $j = 1, \dots, n - 1$, and that $\sigma_1 < 1$, under the assumption that $e^\varphi h_\refe^{-1}$ belongs to $W^{1,2}_{loc}$. However, since the proofs are essentially the same for all cases, we will, for simplicity, only present the proof for $\sigma_1$. We use the mean value inequality (see Theorem \ref{MVI} in Section \ref{CYMMVI}). Since we have $e^{\sigma_1}, e^{\sigma_2}\leq 1$, from (\ref{Hsigma}), it holds that
\begin{align}
\inum\Lambda_{H_1}\partial \bar{\partial} (e^{\sigma_1}-1)&\geq e^{\sigma_1}(3e^{\sigma_1}-3) \geq 3(e^{\sigma_1}-1).
\end{align}
We set $u\coloneqq 1-e^{\sigma_1}$ and denote by $g$ the K\"ahler metric induced by $H_1$. Suppose that $e^{\varphi}h_\refe^{-1}$ belongs to $W^{1,2}_{loc}$. Then the function $u$ also belongs to $W^{1,2}_{loc}$. Let $x_0 \in X$ be an arbitrary point. Then from the mean value inequality theorem, there exists a positive constant $C$ and $0<p<1$ such that
\begin{align}
u(x_0)\geq C\left(\int_{B_g(x_0,1/4)}u^pd\mu_g\right)^{1/p}. \label{ux0}
\end{align}
Let $A_u$ be the set $\{x\in X\mid u(x)=0\}$. Inequality (\ref{ux0}) implies that the set $A_u$ is both closed and open. By the same reason as in the proof of Lemma \ref{26}, $u$ is a non-constant function. Therefore, it follows that $A_u$ is an empty set, and thus, we have $e^{\sigma_1}<1$. 
\end{proof}


\begin{lemm}\label{PQ}{\it Let $P_0\leq P_1\leq \cdots \leq P_{r-1}$ and $Q_0\leq Q_1\leq \dots \leq Q_{r-1}$ be non-negative real numbers such that $P_0+\cdots +P_{r-1}=1$, $Q_0+\cdots +Q_{r-1}=1$. Suppose that $P_1,\dots, P_{r-1}$ and $Q_1,\dots, Q_{r-1}$ are strictly positive and that the following holds:
\begin{align*}
Q_j/Q_{j+1}\leq P_j/P_{j+1} \ \text{for all $j=0,\dots, r-1$.}
\end{align*}
Then the Shannon entropy satisfies 
\begin{align*}
S(Q_0, Q_1,\dots, Q_{r-1})\leq S(P_0, P_1,\dots, P_{r-1}).
\end{align*}
Moreover, the equality is achived if and only if $Q_j/Q_{j+1}=P_j/P_{j+1} $ holds for all $j=1,\dots, r-1$.
} 
\end{lemm}
\begin{proof}
We set $\widetilde{S}(t_1,\dots, t_{r-1})\coloneqq S(1-\sum_{j=1}^{r-1}t_j, t_1,\dots, t_{r-1})$ for variables $t_1,\dots, t_{r-1}$ that take values in $\{t\mid 0<t\leq 1\}$. We set $t_0$ as $t_0\coloneqq 1-\sum_{j=1}^{r-1}t_j$ and $s_0, s_1, \dots, s_{r-1}$ as 
\begin{align}
s_j&\coloneqq t_j/t_{j+1} \ \text{for $j=0,1, \dots, r-2$}, \label{stt} \\
s_{r-1}&\coloneqq 1.
\end{align}
We also introduce the symbols $s^{(l)}\ (l=0,1,\dots, r-1)$ as follows:
\begin{align*}
s^{(l)}&\coloneqq \Pi_{j=l}^{r-1}s_j=s_l\cdot s_{l+1}\cdot\cdots s_{r-1} \ \text{for $l=0,\dots, r-1$}.
\end{align*}
By multiplying both sides of (\ref{stt}) for all $j=0,\dots, r-2$, we obtain
\begin{align}
s^{(0)}=t_0/t_{r-1}. \label{s0}
\end{align}
From (\ref{stt}), we have
\begin{align}
t_j=t_{j+1}s_j \ \text{for all $j=0,1,\dots, r-2$}. \label{tts}
\end{align}
From (\ref{tts}), we have
\begin{align}
\sum_{j=1}^{r-1}t_j=\left(\sum_{l=1}^{r-1}s^{(l)}\right)t_{r-1}. \label{sumtj}
\end{align}
Therefore from (\ref{s0}) and (\ref{sumtj}), we have
\begin{align*}
t_{r-1}s^{(0)}=1-\left(\sum_{l=1}^{r-1}s^{(l)}\right)t_{r-1}
\end{align*}
and thus
\begin{align}
t_{r-1}=\frac{1}{\sum_{l=0}^{r-1}s^{(l)}}. \label{tr-1}
\end{align}
From (\ref{stt}) and (\ref{tr-1}), we obtain
\begin{align}
t_j=\frac{s^{(j)}}{\sum_{l=0}^{r-1}s^{(l)}} \ \text{for $j=0,1,\dots, r-1$}.
\end{align}
By a direct calculation, for each $j=1,\dots, r-1$, we have
\begin{align}
\frac{\partial\widetilde{S}}{\partial t_j}=\log(t_0/t_j)=\sum_{k=0}^{j-1}\log(s_k).
\end{align}
Also, for each $j=1,\dots, r-1$ and $k=0,\dots, r-2$, we have
\begin{align}
\frac{\partial t_j}{\partial s_k}&=\frac{1}{(\sum_{l=0}^{r-1}s^{(l)})^2}\left(-s^{(j)}\frac{\partial}{\partial s_k}\left(\sum_{l=0}^{r-1}s^{(l)}\right)+\frac{\partial s^{(j)}}{\partial s_k}\sum_{l=0}^{r-1}s^{(l)}\right). \label{tjsk}
\end{align}
Suppose that $j>k$. Then we have
\begin{align*}
\frac{\partial s^{(j)}}{\partial s_k}=0.
\end{align*}
Thereofore from (\ref{tjsk}), we obtain 
\begin{align*}
\frac{\partial t_j}{\partial s_k}<0. 
\end{align*}
Suppose that $j\leq k$. Then we have
\begin{align*}
\frac{\partial t_j}{\partial s_k}&=\frac{1}{(\sum_{l=0}^{r-1}s^{(l)})^2}\left(-s^{(j)}\frac{\partial}{\partial s_k}\left(\sum_{l=0}^{r-1}s^{(l)}\right)+\frac{\partial s^{(j)}}{\partial s_k}\sum_{l=0}^{r-1}s^{(l)}\right) \\
&=\frac{\frac{\partial s^{(j)}}{\partial s_k}}{(\sum_{l=0}^{r-1}s^{(l)})^2}\left(-s_k\frac{\partial}{\partial s_k}\left(\sum_{l=0}^{r-1}s^{(l)}\right)+\sum_{l=0}^{r-1}s^{(l)}\right) \\
&=\frac{\frac{\partial s^{(j)}}{\partial s_k}}{(\sum_{l=0}^{r-1}s^{(l)})^2}\left(-\sum_{l=0}^ks^{(l)}+\sum_{l=0}^{r-1}s^{(l)}\right) >0.
\end{align*}
From the above calculations, for each $k=0,1,\dots, r-2$, we have
\begin{align}
\frac{\partial S}{\partial s_k}&=\sum_{j=1}^{r-1}\frac{\partial S}{\partial t_j}\frac{\partial t_j}{\partial s_k} \notag \\
&=\sum_{j=1}^{r-1}(\sum_{l=0}^{j-1}\log (s_l))\frac{\partial t_j}{\partial s_k} \notag \\
&=\sum_{l=0}^{r-1}\log (s_l) \sum_{j=l+1}^{r-1}\frac{\partial t_j}{\partial s_k} \notag \\
&=\sum_{l=0}^{r-1}\log (s_l) \frac{\partial}{\partial s_k}\sum_{j=l+1}^{r-1}t_j \notag \\
&=\sum_{l=0}^{k-1}\log (s_l) \frac{\partial}{\partial s_k}\sum_{j=l+1}^{r-1}t_j+\sum_{l=k}^{r-1}\log (s_l) \frac{\partial}{\partial s_k}\sum_{j=l+1}^{r-1}t_j\notag \\
&=\sum_{l=0}^{k-1}\log (s_l) \frac{\partial}{\partial s_k}\sum_{j=0}^l(-t_j)+\sum_{l=k}^{r-1}\log (s_l) \frac{\partial}{\partial s_k}\sum_{j=l+1}^{r-1}t_j>0. \label{congratulations}
\end{align}
From (\ref{congratulations}), we have the claim. 
\end{proof}

From Proposition \ref{r/2} and Lemma \ref{PQ}, we have Theorem \ref{main theorem a}.

\subsection{Proof of Theorem \ref{main theorem b}}
We next prove Theorem \ref{main theorem b}.
\begin{proof}[Proof of Theorem \ref{main theorem b}] 
We set $Z_{r,\beta}$ and $h_{r,\beta}$ as follows:
\begin{align}
&Z_{r,\beta}\coloneqq \sum_{j=1}^{r-1}\lambda_j^\beta, \label{naze}\\
&h_{r,\beta}\coloneqq \sum_{j=1}^{r-1}\frac{1}{r}\left(\frac{\lambda_j^\beta}{r^{2\beta}}\right) \log\left(\frac{\lambda_j^\beta}{r^{2\beta}}\right).
\label{er}
\end{align}
Then $S_{r,\beta}$ is calculated as follows:
\begin{align}
S_{r,\beta}&=-\sum_{j=1}^{r-1}\left(\frac{\lambda_j^\beta}{Z_{r,\beta}}\right)\log\left(\frac{\lambda_j^\beta}{Z_{r,\beta}}\right) \notag \\ 
&=-\sum_{j=1}^{r-1}\frac{\lambda_j^\beta}{Z_{r,\beta}} \left(\log\frac{\lambda_j^\beta}{r^{2\beta}}-\log r+\log\frac{r^{2\beta+1}}{Z_{r,\beta}}\right) \notag \\
&=-\frac{r^{2\beta+1}}{Z_{r,\beta}}\sum_{j=1}^{r-1}\frac{1}{r}\left(\frac{\lambda_j^\beta}{r^{2\beta}}\right) \log\left(\frac{\lambda_j^\beta}{r^{2\beta}}\right)+\log r-\log\left(\frac{r^{2\beta+1}}{Z_{r,\beta}}\right) \notag \\
&=-\frac{h_{r,\beta}}{(Z_{r,\beta}/r^{2\beta+1})}+\log r+\log\left(\frac{Z_{r,\beta}}{r^{2\beta+1}}\right). \label{Srb}
\end{align}
We define constants $c_\beta$ and $d_\beta$, which may not be finite, as follows:
\begin{align*}
&c_\beta\coloneqq \int_0^1s^\beta(1-s)^\beta ds, \\
&d_\beta\coloneqq \int_0^1s^\beta(1-s)^\beta \log s\ ds.
\end{align*}
Note that each of the above constants is finite if and only if $\beta\leq -1$, and if $\beta\leq -1$, then $c_\beta=\infty$ and $d_\beta=-\infty$. Suppose that $\beta$ is positive. Then since $Z_{r,\beta}/r^{2\beta+1}$ and $h_{r,\beta}$ are Riemann sums of $f_\beta(s)\coloneqq s^\beta(1-s)^\beta$ and $g_\beta (s)\coloneqq s^\beta(1-s)^\beta\log s^\beta(1-s)^\beta$ on the interval $0\leq s\leq 1$, respectively, we have 
\begin{align*}
\lim_{r\to\infty} (S_{r,\beta}-\log r)=-\frac{2\beta d_\beta}{c_\beta}+\log(c_\beta).
\end{align*}
We next consider the case where $\beta$ is negative. The following estimate holds:
\begin{align}
Z_{r,\beta} /r^{2\beta+1}
&\geq \int_{1/r}^{n/r}s^\beta(1-s)^\beta ds+\int_{n/r}^{1-1/r}s^\beta(1-s)^\beta ds \notag \\
&=\int_{1/r}^{1-1/r}s^\beta(1-s)^\beta ds, \label{1-s}\\
Z_{r,\beta}/r^{2\beta+1}&\leq \int_{1/r}^{n/r}s^\beta (1-s)^\beta ds +\frac{1}{r^{2\beta+1}}(r-1)^\beta+\int_n^{1-1/r}s^\beta (1-s)^\beta ds+\frac{1}{r^{2\beta+1}}(r-1)^\beta \notag\\
&=\int_{1/r}^{1-1/r}s^\beta (1-s)^\beta ds+\frac{2}{r^{2\beta+1}}(r-1)^\beta \label{int}. 
\end{align}
As a consequence of the above, for the case where $\beta\leq -1$, as $r$ goes to $\infty$, we have
\begin{align*}
Z_{r,\beta} /r^{2\beta+1}=\begin{cases}
C_{-1}\log r+o(\log r) \ &\text{if $\beta=-1$}\\
C_\beta r^{-\beta-1}+o(r^{-\beta-1}) \ &\text{if $\beta< -1$},
\end{cases}
\end{align*}
where $C_{-1}$ and $C_\beta$ are positive constants and where we denote by $o(\log r)$ (resp. $o(r^{-\beta-1})$) the terms that diverge strictly slower than $\log r$ (resp. $r^{-\beta-1}$) in the limit $r\to\infty$.
Also, by the monotonicity of the function $g_\beta$ at both ends of the interval, there exists $\epsilon_\beta>0$ such that on the interval $0<s\leq \epsilon_\beta$ (resp. $1-\epsilon_\beta\leq s<1$), $g_\beta$ monotonically decreases (resp. monotonically increases). Therefore, if we take $r$ large enough and $j$ (resp. $k$) small enough (resp. close enough to $r$) so that $j/r\leq \epsilon_\beta$ (resp. $1-\epsilon_\beta\leq k/r$) holds, the following holds:
\begin{align}
&\int_{(j-1)/r}^{j/r}g_\beta(s)ds\geq \frac{1}{r}g_\beta(j/r)\geq \int_{j/r}^{(j+1)/r}g_\beta(s)ds, \label{jr}\\ 
&\int_{(k-1)/r}^{k/r}g_\beta(s)ds\leq \frac{1}{r}g_\beta(k/r)\leq \int_{k/r}^{(k+1)/r}g_\beta(s)ds. \label{kr}
\end{align}
We take $r$ large enough and set $j_0\coloneqq [r\epsilon_\beta ]+1$ and $k_0\coloneqq [r-r\epsilon_\beta]-1$, where we denote by $[\cdot]$ the Gaussian symbol. Note that we have $(j_0-1)/r\leq \epsilon_\beta$ and $(k_0+1)/r\leq 1-\epsilon_\beta$. Then we have the following estimate:
\begin{align}
&h_{r,\beta}\leq \frac{1}{r}g_\beta(1/r)+\int_{1/r}^{(j_0-1)/r}g_\beta(s)ds+\sum_{j=j_0}^{k_0}\frac{1}{r}g_\beta(j/r)+\int_{(k_0+1)/r}^{1-1/r}g_\beta(s)ds+\frac{1}{r}g_\beta(1-1/r), \label{grb}\\
&h_{r,\beta}\geq \int_{1/r}^{j_0/r}g_\beta(s)ds+\sum_{j=j_0}^{k_0}\frac{1}{r}g_\beta(j/r)+\int_{k_0/r}^{1-1/r}g_\beta(s)ds. \label{gbs}
\end{align}
From (\ref{grb}) and (\ref{gbs}), for the case where $\beta>-1$, we have
\begin{align}
\lim_{r\to\infty} h_{r,\beta}&=\int_0^{\epsilon_\beta} g_\beta(s)ds+\int_{\epsilon_\beta}^{1-\epsilon_\beta}g_\beta(s)ds +\int_{\epsilon_\beta}^{1-\epsilon_\beta}g_\beta(s)ds, \notag \\
&=\int_0^1g_\beta(s)ds \noindent \\
&=2\beta\int_0^1s^\beta(1-s^\beta)\log sds.\label{limr}
\end{align}
where we have used the easily verifiable fact that for the case where $\beta>-1$, the terms $\frac{1}{r}g_\beta(1/r)=\frac{1}{r}g_\beta(1-1/r)$ goes to $0$ as $r\to\infty$.
Also, from (\ref{grb}) and (\ref{gbs}), for the case where $\beta\leq -1$, we have
\begin{align}
h_{r,\beta}=\begin{cases}
C_{-1}^\prime(\log r)^2+o((\log r)^2) \ &\text{if $\beta=-1$}\\
C_\beta^\prime r^{-\beta-1}\log r+o(r^{-\beta-1}\log r) \ &\text{if $\beta< -1$},
\end{cases} \label{hrbeta}
\end{align}
where $C^\prime_{-1}$ and $C_\beta^\prime$ are positive constants and where we denote by $o((\log r)^2)$ (resp. $o(r^{-\beta-1}\log r)$) the terms that decay strictly slower than $(\log r)^2$ (resp. $r^{-\beta-1}\log r$) as $r$ goes to $\infty$.
Then we prove the assertion for the case where $\beta$ is negative. Suppose that $\beta>-1$. From (\ref{Srb}), (\ref{1-s}), (\ref{int}), and (\ref{limr}), we have 
\begin{align*}
\lim_{r\to\infty} (S_{r,\beta}-\log r)=-\frac{2\beta d_\beta}{c_\beta}+\log(c_\beta).
\end{align*}
We next suppose that $\beta\leq -1$. Then from (\ref{Srb}), (\ref{1-s}), (\ref{int}), and (\ref{hrbeta}), we have
\begin{align*}
S_{r,\beta}-\log r=
\begin{cases}
-\frac{C_{-1}^\prime(\log r)^2+o((\log r)^2)}{C_{-1}\log r+o(\log r)}+\log\left(C_{-1}\log r+o(\log r)\right) \ &\text{if $\beta=-1$}\\
-\frac{C_\beta^\prime r^{-\beta-1}\log r+o(r^{-\beta-1}\log r)}{C_\beta r^{-\beta-1}+o(r^{-\beta-1})}+\log\left(C_\beta r^{-\beta-1}+o(r^{-\beta-1})\right) \ &\text{if $\beta< -1$}.
\end{cases}
\end{align*}
For both cases, we have $\lim_{r\to\infty}(S_{r,\beta}-\log r)=-\infty$. This implies the claim.
\end{proof}




\begin{thebibliography}{99}
\bibitem{AF1}E. Aldrovandi and G. Falqui, {\it Geometry of Higgs and Toda fields on Riemann surfaces}, J. Geom. Phys. 17 (1995), no. 1, 25-48.
\bibitem{Bar1} D. Baraglia, {\it $G_2$ geometry and integrable systems}, arXiv:1002.1767 ver2 (2010).
\bibitem{BCHM1} T. Bayraktar, D. Coman, H. Herrmann, and G. Marinescu, {\it A survey on zeros of random holomorphic sections}, arXiv:1807.05390 ver2 (2018).
\bibitem{BH1} Y. Benoist and D. Hulin, {\it Cubic differentials and hyperbolic convex sets}, Journal of Differential Geometry 98.1 (2014): 1-19.
\bibitem{BB1} R. Berman and S. Boucksom, {\it Growth of balls of holomorphic sections and energy at equilibrium} Inventiones mathematicae, 181(2), 337-394 (2010).
\bibitem{BS1} T. Bridgeland and I. Smith, {\it Quadratic differentials as stability conditions}, Publications math\'ematiques de l'IH\'ES 121 (2015): 155-278.
\bibitem{CY1} S. Y. Cheng and S. T. Yau, {\it Differential equations on Riemannian manifolds and their geometric applications}, Communications on Pure and Applied Mathematics 28.3 (1975): 333-354.
\bibitem{CT1} H. I. Choi and A. Treibergs, {\it Gauss maps of spacelike constant mean curvature hypersurfaces of Minkowski space}, Journal of differential geometry 32.3 (1990): 775-817.
\bibitem{Col1} B. Collier, {\it Finite order automorphisms of Higgs bundles: theory and application}, (Doctoral dissertation, University of Illinois at Urbana-Champaign).
\bibitem{CL1} B. Collier and Q. Li, {\it Asymptotics of Higgs bundles in the Hitchin component}, Advances in Mathematics 307 (2017): 488-558.
\bibitem{DL1} S. Dai and Q. Li, {\it Minimal surfaces for Hitchin representations}, Journal of Differential Geometry 112.1 (2019): 47-77.
\bibitem{DL2} S. Dai and Q. Li, {\it On cyclic Higgs bundles}, Math. Ann. 376 (2020), no. 3-4, 1225-1260.
\bibitem{DL3} S. Dai and Q. Li, {\it Bounded differentials on the unit disk and the associated geometry}, Transactions of the American Mathematical Society (2024).
\bibitem{Don1} S. K. Donaldson, {\it Boundary value problems for Yang-Mills fields}, Journal of geometry and physics 8.1-4 (1992): 89-122.
\bibitem{DW1} D. Dumas and M. Wolf, {\it Polynomial cubic differentials and convex polygons in the projective plane}, Geometric and Functional Analysis 25.6 (2015): 1734-1798.
\bibitem{FKMR1} O. El-Fallah, K. Kellay, J. Mashreghi, and T. Ransford, {\it A primer on the Dirichlet space} (Vol. 203). Cambridge University Press (2014).
\bibitem{GZ1} V. Guedj and A. Zeriahi, {\it Degenerate Complex Monge-Amp\`ere Equations}, EMS Tracts Math. 26, European Mathematical Society (EMS), Z\"urich 2017.
\bibitem{GL1} M. A. Guest and C.-S. Lin, {\it Nonlinear PDE aspects of the $tt^\ast$ equations of Cecotti and Vafa}, J. Reine Angew. Math. 689 (2014), 1-32.
\bibitem{GH1} M. A. Guest and N. K. Ho, {\it Kostant, Steinberg, and the Stokes matrices of the $tt^\ast$-Toda equations}, Selecta Mathematica 25.3 (2019): 50.
\bibitem{Hit1} N. J. Hitchin, {\it The self-duality equations on a Riemann surface}, Proceedings of the London Mathematical Society 3.1 (1987): 59-126.
\bibitem{Hit2} N. J. Hitchin, {\it Lie groups and Teichm\"uller space}, Topology, 31(3), 449-473 (1992).
\bibitem{JT1} A. Jaffe and C. Taubes, {\it Vortices and monopoles}, vol. 2 of Progress in Physics (1980).
\bibitem{KH1} A. Katok and B. Hasselblatt, {\it Introduction to the modern theory of dynamical systems (No. 54)}, Cambridge university press (1995). 
\bibitem{Kus1} Y. Kusunoki, {\it Integrable holomorphic quadratic differentials with simple zeros}, Holomorphic Functions and Moduli I: Proceedings of a Workshop held March 13-19, 1986. Springer US, 1988.
\bibitem{LT1} F. Labourie and J. Toulisse, {\it Quasicircles and quasiperiodic surfaces in pseudo-hyperbolic spaces}, Inventiones mathematicae 233.1 (2023): 81-168.
\bibitem{LTW1} F. Labourie, J. Toulisse, and M. Wolf, {\it Plateau problems for maximal surfaces in pseudo-hyperbolic spaces}, arXiv:2006.12190 ver 3 (2020).
\bibitem{LL1} L. D. Landau and E. M. Lifshitz, {\it Statistical Physics: Volume 5}, Vol. 5. Elsevier, 2013.
\bibitem{Lei1} T. Leinster, {\it Entropy and Diversity: The Axiomatic Approach}, Cambridge University Press (2020).
\bibitem{Li0} Q. Li, {\it Harmonic maps for Hitchin representations}, Geometric and Functional Analysis 29 (2019): 539-560.
\bibitem{Li1} Q. Li, {\it On the uniqueness of vortex equations and its geometric applications}, The Journal of Geometric Analysis 29 (2019): 105-120.
\bibitem{LM1} Q. Li and T. Mochizuki, {\it Complete Solutions of Toda Equations and Cyclic Higgs Bundles Over Non-compact Surfaces}, International Mathematics Research Notices 2025.7 (2025): rnaf081.
\bibitem{LM2} Q. Li and T. Mochizuki, {\it Isolated singularities of Toda equations and cyclic Higgs bundles}, arXiv:2010.06129 ver 3 (2020).
\bibitem{LM3} Q. Li and T. Mochizuki, {\it Higgs bundles in the Hitchin section over non‐compact hyperbolic surfaces}, Proceedings of the London Mathematical Society 129.6 (2024): e70008.
\bibitem{Lof1} J. C. Loftin, {\it Flat metrics, cubic differentials and limits of projective holonomies}, Geometriae Dedicata, 128 (2007), 97-106.
\bibitem{LoTW1} J. Loftin, A. Tamburelli, and M. Wolf., {\it Limits of cubic differentials and buildings}, arXiv:2208.07532 (2022).
\bibitem{Ma1} W. Ma, {\it Harmonic metrics of $\mathrm {SO} _ {0}(n, n) $-Higgs bundles in the Hitchin section on non-compact hyperbolic surfaces}, arXiv:2408.15278 ver 1 (2024).
\bibitem{Mas1} H. Masur, {\it The growth rate of trajectories of a quadratic differential}, Ergodic Theory and Dynamical Systems 10.1 (1990): 151-176.
\bibitem{Miy2} N. Miyatake, {\it Generalizations of Hermitian-Einstein equation of cyclic Higgs bundles, their heat equation, and inequality estimates}, arXiv:2301.01584 ver2 (2023).
\bibitem{Miy3} N. Miyatake, {\it Cyclic Higgs bundles, subharmonic functions, and the Dirichlet problem}, Annals of Global Analysis and Geometry 67.1 (2025): 1-16.
\bibitem{Miy4} N. Miyatake, {\it Shannon entropy for harmonic metrics on cyclic Higgs bundles II}, preprint.
\bibitem{Miy5} N. Miyatake, {\it Complete harmonic metrics and subharmonic functions on the unit disc}, preprint. 
\bibitem{Moc0} T. Mochizuki,{\it Harmonic bundles and Toda lattices with opposite sign I}, arXiv:1301.1718.
\bibitem{Moc1} T. Mochizuki, {\it Harmonic bundles and Toda lattices with opposite sign II.}, Communications in Mathematical Physics 328 (2014): 1159-1198.
\bibitem{Nie1} X. Nie, {\it Poles of cubic differentials and ends of convex $\mathbb {RP}^ 2$-surfaces}, Journal of Differential Geometry 123.1 (2023): 67-140.
\bibitem{PW1} Y. Polyanskiy and Y. Wu, {\it Lecture notes on information theory}, Lecture Notes for ECE563 (UIUC) and 6.2012-2016 (2014): 7.
\bibitem{Ran1} T. Ransford, {\it Potential theory in the complex plane}, No. 28. Cambridge university press, 1995.
\bibitem{ST1} E. B. Saff and V. Totik, {\it Logarithmic potentials with external fields}, Vol. 316. Springer Science \& Business Media, 2013.
\bibitem{Sak1} K. Sakai, {\it Uniform degeneration of hyperbolic surfaces with boundary along harmonic map rays}, arXiv:2405.0485 ver2 (2024).
\bibitem{Ser1} S. Serfaty, {\it Coulomb gases and Ginzburg-Landau vortices}, arXiv:1403.6860 ver 4 (2014).
\bibitem{Ser2} S. Serfaty, {\it Systems with Coulomb interactions}, Journ\'ees \'equations aux d\'eriv\'ees partielles (2014): 1-23.
\bibitem{Sh1} C. E. Shannon, {\it A mathematical theory of communication}, The Bell system technical journal, 27(3), 379-423 (1948).
\bibitem{SZ1} B. Shiffman and S. Zelditch, {\it Stochastic K\" ahler geometry: from random zeros to random metrics}, arXiv:2303.11559 ver 1 (2023).
\bibitem{Sim1} C. T. Simpson, {\it Constructing variations of Hodge structure using Yang-Mills theory and applications to uniformization}, Journal of the American Mathematical Society (1988): 867-918.
\bibitem{Sim2} C. T. Simpson, {\it Harmonic bundles on noncompact curves}, Journal of the American Mathematical Society 3.3 (1990): 713-770.
\bibitem{Sim3} C. T. Simpson, {\it Katz's middle convolution algorithm}, arXiv preprint math/0610526. ver2 (2006).
\bibitem{St1} K. Strebel, {\it Quadratic differentials}, Springer Berlin Heidelberg, 1984.
\bibitem{Wan1} T. Y.-H. Wan, {\it Constant mean curvature surface, harmonic maps, and universal Teichm\"uller space}, J. Differential Geom.35 (1992), no.3, 643-657.
\bibitem{WA1} T. Y.-H. Wan and T. K.-K. Au, {\it Parabolic constant mean curvature spacelike surfaces}, Proceedings of the American Mathematical Society 120.2 (1994): 559-564.
\bibitem{Wol1} M. Wolf, {\it The Teichm\"uller theory of harmonic maps}, Journal of differential geometry 29.2 (1989): 449-479.
\end{thebibliography}
\noindent
E-mail address 1: natsuo.miyatake.e8@tohoku.ac.jp

\noindent
E-mail address 2: natsuo.m.math@gmail.com \\

\noindent
Mathematical Science Center for Co-creative Society, Tohoku University, 468-1 Aramaki Azaaoba, Aoba-ku, Sendai 980-0845, Japan.



\end{document}


























